% ============================================================
% Biomedical Signal Processing and Control (Elsevier)
% Tri-Modal Depression Risk Detection
% Author: Md. Mursalin, Netrokona University
% ============================================================

\documentclass[5p,twocolumn]{elsarticle}

% ---- Packages ----
\usepackage{amsmath,amssymb,amsfonts}
\usepackage{algorithm}
\usepackage{algorithmic}
\usepackage{graphicx}
\usepackage{textcomp}
\usepackage{xcolor}
\usepackage{booktabs}
\usepackage{multirow}
\usepackage{url}
\usepackage[hidelinks]{hyperref}
\journal{Biomedical Signal Processing and Control}

% In Overleaf: upload all fig*.png to project root
\graphicspath{{./}}

\begin{document}

\begin{frontmatter}

\title{A Systematic Validity Audit of Tri-Modal Depression Detection on
DAIC-WOZ: Sentinel Artifacts, Interviewer Shortcuts, and Cross-Corpus
Transfer}

\author[netrokona]{Md.~Mursalin\corref{cor1}\fnref{code}}
\ead{mursalinshuvo27@gmail.com}
\author[netrokona]{Md.~Shovon}

\cortext[cor1]{Corresponding author.}
\fntext[code]{Dataset: DAIC-WOZ provided by USC Institute for
Creative Technologies. Code:
\url{https://github.com/DevMursLab/Thesis} (archived version, DOI:
\url{https://doi.org/10.5281/zenodo.22806159}).}
\address[netrokona]{Department of Computer Science and Engineering,
Netrokona University, Netrokona 2400, Bangladesh}

% ============================================================
\begin{abstract}
Reported performance for automated depression screening on the
DAIC-WOZ corpus varies by more than 20 AUC points across the
literature under seemingly minor pipeline differences. We conduct a
systematic validity audit of a tri-modal (facial action units,
acoustic features, transcript text) depression-risk pipeline and
surface three failure modes that plausibly account for a substantial
share of this cross-study variance: (1)~a $-100$ missing-data
sentinel in OpenFace's binary AU channels that survives standard
confidence filtering and, under per-participant z-scoring, compresses
the genuine facial signal $\approx$79$\times$, holding a face-only
model at chance; removing sentinel-carrying frames raises face-only
AUC from $0.506$ to $0.678$ ($n{=}30$ paired seeds, 28/30 improving,
Wilcoxon $p{=}9.3{\times}10^{-6}$, $d{=}1.28$); (2)~a direct
demonstration that DAIC-WOZ text models can learn interview structure
instead of depression language---an interviewer-only model with zero
access to patient words reaches test AUC~$=0.738$, exceeding our
participant-only model's AUC~$=0.610$; and (3)~cross-corpus transfer
collapsing to chance for both audio (EATD-Corpus, AUC~$0.735\to0.531$)
and text (WU3D, 32,540 users, AUC~$0.506$), indicating the learned
representations are corpus/language-specific rather than transferable
biomarkers. Within the same controlled testbed, symptom-level
multi-task learning raises mean macro-F1 from $0.607$ to $0.629$ at
an originally-reported, underpowered 5-seed sample ($n{\approx}67$
needed for 80\% power); scaling to $n{=}70$ paired seeds confirms the
effect of the combined intervention (multi-task heads plus extended
audio features; macro-F1 $0.597\pm0.076\to0.657\pm0.037$, Wilcoxon
$p{=}8.1{\times}10^{-7}$, $d{=}0.72$), while cross-modal attention and
an Equalized Odds fairness loss remain non-significant at $n{=}10$
seeds and are reported as transparent negative/inconclusive findings.
On the held-out test split ($N{=}47$), the model achieves
F1~$=0.585$, AUC~$=0.706$, exceeding dev AUC and indicating
generalization; pooled 10-fold cross-validation over all 188
participants shows the test split's nominal gender-TPR parity is not
stable across folds (mean TPR gap~$0.415\pm0.211$). Three convergent
explainability methods identify audio as the most informative
modality on dev, and gradient saliency places the clinically-known
marker AU04 within its top-5 attributed facial features, though not
top-ranked. An external diagnostic-specificity probe finds the model
separates depression from normal controls (AUC~$0.678$) but only
weakly from suicidality ($0.573$) or anxiety ($0.568$), partially
supporting the concern that DAIC-WOZ classifiers capture
disorder-general rather than depression-specific signal. We distill
these audits into a six-item, reusable validity checklist for other
DAIC-WOZ studies. Our contribution is methodological rigor and
transparent reporting---including honest negative results and
explicit statistical-power analysis---rather than
state-of-the-art performance.
\end{abstract}

\begin{keyword}
Depression detection \sep multi-task learning \sep PHQ-8 \sep
multimodal fusion \sep DAIC-WOZ \sep CLNF action units \sep
cross-modal attention \sep fairness-aware training \sep
Equalized Odds \sep explainable AI
\end{keyword}

\end{frontmatter}

% ============================================================
\section{Introduction}
\label{sec:intro}

Depression affects an estimated 280 million people
worldwide \cite{who2023} and is a leading cause of disability globally.
Clinical diagnosis relies on structured interviews, self-report
questionnaires (PHQ-8, PHQ-9), and clinician observation---a
resource-intensive process that is inaccessible in many healthcare
settings. Automated screening from clinical interview recordings offers
a scalable complement to traditional pathways, provided the system
meets three requirements: (1)~\emph{multimodal coverage} matching the
multi-channel nature of clinical observation; (2)~\emph{demographic
fairness} ensuring equal detection rates across gender groups; and
(3)~\emph{clinical interpretability} enabling a human clinician to
audit the model's reasoning.

Prior work on the DAIC-WOZ corpus \cite{gratch2014} has explored
audio features \cite{ringeval2017}, text-based models, and multimodal
fusion \cite{gong2017,williamson2016}. However, three critical gaps
persist. First, ablation studies in existing work typically compare
configurations with \emph{different} architectural choices per modality
subset, confounding the modality contribution with incidental model
changes. Second, fairness constraints---if present---address only
Equal Opportunity (TPR gap), which can be satisfied trivially by
predicting positive for all participants. Third, explainability
evidence is typically provided by a single method without
cross-validation.

DAIC-WOZ is the de facto benchmark for automated depression screening,
yet reported results across the literature vary by 20+ AUC points
under seemingly minor pipeline differences, and few studies test
whether their models measure depression at all as opposed to
interview structure, tracking-software artifacts, or corpus-specific
acoustic idiosyncrasies. This paper's central contribution is a
systematic validity audit that surfaces and directly tests three
such failure modes---each independently capable of explaining a
substantial share of the variance in reported DAIC-WOZ performance
across the field---using a controlled tri-modal ablation testbed
built specifically to isolate them from ordinary architectural
confounds. Concretely, we contribute:

\begin{enumerate}
  \item \textbf{Three validity-failure-mode audits, each substantiated
  by a controlled test in our own pipeline.} (a)~\emph{A missing-data
  sentinel that silently disables an entire modality:} a $-100$
  tracking-failure sentinel in OpenFace's binary AU channels survives
  the standard confidence filter and, under per-participant z-scoring,
  compresses the genuine signal $\approx$79$\times$
  (Section~\ref{sec:sentinel}); removing the affected frames raises
  face-only AUC from $0.506$ to $0.678$ across $n{=}30$ paired seeds
  (28/30 seeds improving, Wilcoxon $p{=}9.3{\times}10^{-6}$,
  $d{=}1.28$)---this is
  the kind of silent data-corruption bug that can produce a false null
  result in any pipeline reusing OpenFace output without auditing
  sentinel values. (b)~\emph{DAIC-WOZ text models can learn interview
  structure instead of depression language:} building on a concern
  raised by recent reproducibility work
  \cite{burdisso2024therapist,patapati2025invalid}, an interviewer-only
  text model---with zero access to any word the patient
  says---reaches test AUC~$=0.738$, \emph{exceeding} our
  participant-only model's AUC~$=0.610$ (Section~\ref{sec:shortcut}).
  (c)~\emph{Cross-corpus transfer collapses to chance for both audio
  and text:} our strongest unimodal (audio) model collapses from
  in-domain AUC~$=0.735$ to AUC~$=0.531$ (95\% CI $[0.406,0.653]$) on
  EATD-Corpus (Chinese-language); an independent text-modality probe
  on WU3D \cite{wang2020wu3d} (32,540 Weibo users) replicates this with
  a far tighter bound (AUC~$=0.506$, 95\% CI $[0.5008,0.5121]$), traced
  to 55.9\% of cross-lingual documents producing all-zero TF-IDF
  vectors. All three indicate failure modes that plausibly generalize
  beyond our own pipeline to other DAIC-WOZ studies that do not audit
  for them.

  \item \textbf{A controlled tri-modal ablation testbed, and a
  reusable validity checklist distilled from the audits above.} A
  single \texttt{ConfigurableFusionModel} instantiates all seven
  modality subsets with identical encoders, attention, and classifier,
  attributing performance differences solely to modalities rather than
  incidental architectural choices; this is also what makes audit~(a)
  above possible, since a chance-level face-only score cannot be
  dismissed as an architectural weakness of that branch. Within this
  testbed we evaluate symptom-level multi-task learning, cross-modal
  attention, and Equalized Odds fairness training. A power analysis of
  the original 5-seed multi-task test showed it was structurally too
  small for an exact Wilcoxon test to reach significance under any
  outcome ($n{\approx}67$ needed for 80\% power), so we scaled that
  comparison to 70~paired seeds: multi-task learning is a confirmed
  effect (51/70~seeds, Wilcoxon $p{=}8.1{\times}10^{-7}$, $d{=}0.72$),
  a larger effect than the original power calculation assumed.
  Cross-modal attention and Equalized Odds fairness training remain
  evaluated at 10~seeds and do not reach significance at this sample
  size (attention $p{=}0.496$; fairness $p{=}0.813$), with a post-hoc
  power analysis showing the study is genuinely underpowered for these
  two effect sizes rather than demonstrating true null effects
  (Section~\ref{sec:ablation}); we report both as transparent
  negative/inconclusive findings. We distill the three audits above
  into a six-item checklist (Table~\ref{tab:checklist},
  Section~\ref{sec:checklist}) covering sentinel auditing,
  interviewer-text isolation, cross-corpus transfer, split
  disjointness, seed-count power, and diagnostic specificity, so that
  other DAIC-WOZ studies can apply these checks directly rather than
  repeating our post-hoc discovery process.

  \item \textbf{Falsification checks confirming the headline results
  are not artefacts of the sentinel fault or of dev-split modality
  ranking.} Re-training the proposed multi-task model (5 seeds each
  arm) under the corrected face features from audit~(a) changes AUC by
  $-0.009$ ($p{=}0.125$) and F1 by $+0.030$ ($p{=}0.500$), neither
  significant; a complementary analysis of 85~dynamics features on
  raw, unnormalized AUs finds no blunted-affect signal surviving FDR
  correction, isolating the sentinel rather than the normalization
  scheme as the cause. Separately, test-time zero-masking
  (Section~\ref{sec:robustness}) shows the model tolerates losing face
  with negligible cost but loses most discriminative power when text
  is removed (test AUC $0.706\to0.435$)---a partial, modality-dependent
  robustness result---and Section~\ref{sec:sentinel} confirms the
  face channel's dispensability here is genuine redundancy with audio
  and text, not an artefact of the sentinel fault.

  \item \textbf{Convergent explainability and an external
  diagnostic-specificity probe.} Attention rollout,
  gradient$\times$input saliency, and occlusion ablation, applied
  independently to the same predictions, provide stronger attribution
  evidence than any single method alone, though the dev-split
  attribution ranking does not fully predict test-time modality
  necessity (Section~\ref{sec:robustness}). Because DAIC-WOZ's single
  self-report label cannot distinguish a depression detector from a
  general-distress detector, we test this externally on a
  multi-category mental-health corpus with length-matched contrasts:
  the model transfers (depression vs.\ normal AUC~$=0.678$) but
  separates depression from suicidality ($0.573$) and anxiety
  ($0.568$) only weakly---partial support for the disorder-general
  critique \cite{patapati2025invalid}, and a stated safety bound on
  where such a model may be applied.
\end{enumerate}

The remainder of the paper is organized as follows.
Section~\ref{sec:related} reviews related work.
Section~\ref{sec:dataset} describes the DAIC-WOZ dataset.
Section~\ref{sec:features} details feature extraction.
Section~\ref{sec:architecture} presents the model architecture.
Section~\ref{sec:training} describes the training objective.
Sections~\ref{sec:ablation}--\ref{sec:xai} report experimental results.
Section~\ref{sec:discussion} audits two feature-level validity
concerns---a missing-data sentinel in the visual stream
(Section~\ref{sec:sentinel}) and an interviewer-prompt shortcut in the
text stream (Section~\ref{sec:shortcut})---before discussing limitations
and clinical significance. Section~\ref{sec:conclusion} concludes.

% ============================================================
\section{Related Work}
\label{sec:related}

\subsection{Automated Depression Screening on DAIC-WOZ}

The AVEC 2017 challenge \cite{ringeval2017} established baseline F1
scores of 0.50 (audio) and 0.49 (text) on the DAIC-WOZ dev split,
serving as the primary reference point for subsequent work.
Williamson et al.\ \cite{williamson2016} achieved F1~=~0.57 with
multimodal audio-visual handcrafted features. Gong and Poellabauer
\cite{gong2017} combined LDA topic-model text features with acoustic
features to reach F1~=~0.70. Ray et al.\ \cite{ray2019} employed
BERT for transcript encoding alongside audio features, achieving
AUC~=~0.80. Several
of these higher numbers are obtained with interviewer text included;
recent validity audits \cite{burdisso2024therapist,patapati2025invalid}
show this can inflate performance through interview-structure
shortcuts rather than depression-specific language, a concern we
test directly in our own pipeline (Section~\ref{sec:shortcut}).

To our knowledge, no prior DAIC-WOZ study combines all of: (1)~extended
acoustic features (MFCC+COVAREP+FORMANT), (2)~an Equalized Odds
training loss, (3)~symptom-level PHQ-8 multi-task prediction, and
(4)~controlled multi-seed ablation with an identical model class, in a
single end-to-end system. We note that combining these components is
a methodological first; as reported in Sections~\ref{sec:ablation}
and~\ref{sec:fairness}, the Equalized Odds loss did not yield a
statistically significant effect at this sample size, and we report
that outcome honestly rather than as a validated benefit.

More importantly, this controlled testbed is what makes the three
validity audits in Sections~\ref{sec:sentinel}--\ref{sec:shortcut} and
Section~\ref{sec:xai} possible in the first place: because every
modality subset shares an identical encoder, attention mechanism, and
classifier, a chance-level face-only score cannot be dismissed as an
architectural weakness of that branch, forcing the investigation that
uncovered the OpenFace sentinel fault
(Section~\ref{sec:sentinel}). Similarly, holding the model class fixed
while swapping only which text is visible to it isolates the
interviewer-prompt shortcut (Section~\ref{sec:shortcut}) from any
confound introduced by a different text architecture. To our
knowledge, none of the prior DAIC-WOZ studies cited above audit for
either failure mode, and the wide spread of reported AUC/F1 values
across this literature (0.49--0.80) is consistent with at least some
of that variance reflecting uncontrolled versions of the same issues
rather than genuine modelling gains.

\subsection{A Validity Checklist for DAIC-WOZ Studies}
\label{sec:checklist}

The three failure modes audited in this paper (Sections~\ref{sec:sentinel}--\ref{sec:shortcut}
and~\ref{sec:xai}) are not idiosyncrasies of our own pipeline: each
arises from a property of the DAIC-WOZ release itself (an
undocumented OpenFace sentinel value, an interviewer transcript
channel included by default, a single-corpus/single-language design)
that any pipeline built on this corpus can inherit silently. Because
none of the prior work surveyed above reports auditing for these
specific failure modes (Table~\ref{tab:checklist}), we distill our
audits into a checklist so that other DAIC-WOZ studies can check for
each condition directly, at the cost of a few lines of code, rather
than repeating our post-hoc discovery process. Each row names a
concrete, mechanical check; passing it does not guarantee validity,
but failing it identifies a specific, previously undocumented way a
reported result on this corpus can be inflated or deflated.

\begin{table}[t]
\centering
\renewcommand{\arraystretch}{1.3}
\caption{A Validity Checklist for DAIC-WOZ Depression-Detection Studies}
\label{tab:checklist}
\begin{tabular}{p{0.22\columnwidth}p{0.68\columnwidth}}
\toprule
\textbf{Check} & \textbf{What to verify} \\
\midrule
\textbf{1. OpenFace sentinel audit} &
Before z-scoring or otherwise normalizing OpenFace AU channels, confirm
no frame carries the $-100$ tracking-failure sentinel (or any
out-of-range placeholder) past the standard confidence filter. A
sentinel that survives normalization can silently compress genuine
signal by an order of magnitude (Section~\ref{sec:sentinel}) and
produce a false null result for the visual modality without any
error or warning. \\[2pt]
\textbf{2. Interviewer-text isolation} &
Confirm the text pipeline uses \emph{participant-only} utterances, and
if interviewer turns are included, separately report an
interviewer-only ablation (Section~\ref{sec:shortcut}). A text model
that scores well while including interviewer turns may be exploiting
which questions get asked rather than what the patient says. \\[2pt]
\textbf{3. Cross-corpus / cross-lingual transfer check} &
Before describing a learned representation as a depression
\emph{biomarker} rather than a corpus-fit artifact, evaluate it
zero-shot on at least one independent corpus
(Section~\ref{subsec:limitations}, cross-corpus probes); in-corpus
performance alone cannot distinguish the two. \\[2pt]
\textbf{4. Participant-level split disjointness} &
Verify programmatically---not by construction assumption---that no
participant ID appears in more than one of train/dev/test or across
resampling folds. Subject-level leakage is reported as the most
common and most performance-inflating flaw across a 66-paper
DAIC-WOZ reproducibility audit \cite{danylenko2026}. \\[2pt]
\textbf{5. Seed-count power check} &
Before reporting a paired seed-level significance test (e.g.\
Wilcoxon) as confirmatory or as a null result, compute the seed count
required for adequate power at the observed effect size. An
underpowered test is uninformative in either direction
(Section~\ref{sec:ablation}); we found this directly by discovering our own
$n{=}5$ multi-task comparison could not reach significance under any
outcome, then resolving it by scaling to $n{=}70$ seeds. \\[2pt]
\textbf{6. Diagnostic-specificity check} &
Because PHQ-8 alone cannot distinguish a depression-specific detector
from a general-distress detector, evaluate the model against at least
one other distress category (e.g.\ anxiety, suicidality) on an
external corpus before describing detections as depression-specific
(Section~\ref{subsec:limitations}, diagnostic-specificity probe). \\
\bottomrule
\end{tabular}
\end{table}

We are not aware of a prior published DAIC-WOZ study that reports
having applied all six checks; applying even a subset before
reporting a headline number would have flagged the specific failure
modes this paper identifies.

\subsection{Multimodal Fusion for Affect}

Attention-based multimodal fusion is well established in sentiment
analysis \cite{tsai2019multimodal} and emotion recognition. The
Multimodal Transformer \cite{tsai2019multimodal} uses cross-modal
attention to align sequences across modalities. However, in clinical
depression screening, most fusion approaches concatenate modality
features \cite{williamson2016} or use late fusion without attention
weights, providing no interpretability signal. Our work applies
cross-modal attention with residual connections and empirically
evaluates its advantage over concatenation under controlled
conditions; we find a positive but not statistically significant
trend at this sample size (Section~\ref{subsec:attn_vs_concat}).

\subsection{Fairness in Clinical AI}

Hardt et al.\ \cite{hardt2016} formalize Equalized Odds as requiring
equal TPR \emph{and} FPR across protected groups. Applied to clinical
screening: (a)~depressed patients of both genders are caught at equal
rates (TPR), and (b)~non-depressed patients experience equal false
alarm rates (FPR). Prior depression detection work has not integrated
Equalized Odds into the training objective---existing fairness audits
are post-hoc rather than adversarial. Barocas et al.\ \cite{barocas2019}
provide the formal framework for our four-criterion audit.

% ============================================================
\section{Dataset}
\label{sec:dataset}

\subsection{DAIC-WOZ}

The Distress Analysis Interview Corpus---Wizard of Oz (DAIC-WOZ)
\cite{gratch2014} contains 188 clinical interview recordings from
USC Institute for Creative Technologies. Participants were interviewed
by an animated virtual agent (Ellie); each session was transcribed
and annotated with PHQ-8 self-report scores. Binary depression label:
PHQ-8~$\geq 10$.

We use the official AVEC 2017 split:
\begin{itemize}
  \item \textbf{Train:} 107 participants (30 depressed, 77 not;
    prevalence $\approx28\%$)
  \item \textbf{Dev:} 34 participants (12 depressed, 22 not)
  \item \textbf{Test:} 47 participants (14 depressed, 33 not)
\end{itemize}

Gender labels are provided in the official split CSVs. PHQ-8 scores
and binary labels for all splits---including test---are available to
registered DAIC-WOZ users under the USC ICT data-use agreement.
Model selection and threshold tuning are performed on the dev split
only; the test split is used exclusively for final held-out evaluation.
The training set class
imbalance ratio (77:30~$\approx$~2.6:1) motivates class-weighted
cross-entropy with an upper cap of $1.8\times$ to prevent over-correction.

Splits are participant-disjoint: each participant ID appears in
exactly one of train, dev, or test, so no interview session is
shared across splits. This matters because a recent large-scale
reproducibility audit of DAIC-WOZ studies found that of 66 published
papers reviewed, only 5 met minimal reproducibility standards, and
identified subject-level data leakage across splits as the most
common and most performance-inflating flaw \cite{danylenko2026}. We
verified the absence of participant overlap programmatically across
train/dev/test and, separately, across the k-fold partitions used
in Section~\ref{sec:kfold} (each fold is drawn by \texttt{StratifiedKFold}
over the pooled 188 participants with no repeated IDs across
train/test within a fold).

% ============================================================
\section{Feature Extraction}
\label{sec:features}

\subsection{Facial Action Units (Face)}

OpenFace CLNF \cite{baltrusaitis2016} extracts 20 Action Unit (AU) channels
per frame: 14~continuous AU intensities ($r$ suffix) and 6~binary
presence indicators ($c$ suffix), including AU04 (Brow Lowerer),
AU15 (Lip Corner Depressor), and AU17 (Chin Raiser)---FACS units
\cite{ekman1978} validated as depression affect markers. Frames with confidence
$<0.5$ are excluded; remaining frames are z-score normalized
per channel and uniformly sampled to $T_f=200$ steps.
Output shape: $\mathbf{X}_\text{face}\in\mathbb{R}^{N\times200\times20}$.

\subsection{Clinical Audio (Audio)}

We combine three complementary acoustic feature sets:

\textbf{MFCC+$\Delta$+$\Delta^2$} (120-dim): 40~Mel-frequency
cepstral coefficients plus first- and second-order delta features,
extracted at 16~kHz with hop~length~512. Delta features capture
temporal speech dynamics; depression correlates with reduced prosodic
variability \cite{ringeval2017}.

\textbf{COVAREP} (74-dim): Per-frame features including fundamental
frequency F0 (pitch), Normalized Amplitude Quotient (NAQ), Quasi-Open
Quotient (QOQ), H1H2 harmonic difference, Peak Slope, Rd glottal
parameter, and 24~Mel Cepstral Prediction (MCEP) coefficients
\cite{covarep2014}. Depression manifests as monotone delivery,
reduced F0 range, and breathier voice quality.

\textbf{FORMANT} (5-dim): F1--F5 formant frequencies capturing
vocal tract resonance. Scherer et al.\ \cite{scherer2016vowel}
demonstrated reduced vowel space (compressed F1-F2 dispersion) in
individuals self-reporting depression/PTSD symptoms, distinct from
the voice-quality features described above \cite{scherer2013voice}.

All three streams are independently z-score normalized, resampled to
$T_a=300$ time steps via linear interpolation, and concatenated along
the feature axis:
$\mathbf{X}_\text{audio}\in\mathbb{R}^{N\times300\times199}$.

\subsection{Transcript Bigrams (Text)}

Participant utterances are extracted from DAIC-WOZ transcripts
(interviewer turns excluded). TF-IDF bigrams with vocabulary
$V=1{,}000$, sublinear TF scaling, and minimum document frequency~2
capture phrase-level depression markers (``feel hopeless'',
``can't sleep''). Vocabulary restriction to 1,000 prevents
overfitting on $N=107$.
Shape: $\mathbf{X}_\text{text}\in\mathbb{R}^{N\times1000}$.

% ============================================================
\section{Model Architecture}
\label{sec:architecture}

The full pipeline is illustrated in Fig.~\ref{fig:architecture}.

\begin{figure}[t]
  \centering
  \includegraphics[width=\columnwidth]{fig1_architecture.png}
  \caption{Tri-modal multi-task architecture. Three modality encoders
    project to a shared 64-dim space; cross-modal attention allows
    each modality to query the other two; three prediction heads
    output binary label, PHQ-8 score, and 8 symptom items.}
  \label{fig:architecture}
\end{figure}

\subsection{Modality Encoders}

Each modality is independently encoded to a shared embedding space
$\mathbb{R}^d$, $d=64$:

\textbf{FaceAUEncoder}: 2-layer bidirectional LSTM, hidden~dim~64,
mean-pooling over time, LayerNorm, GELU, Dropout(0.3). Input:
$\mathbb{R}^{200\times20}$.

\textbf{AudioEncoder}: Identical architecture to FaceAUEncoder,
accepting 199-dimensional input per step.

\textbf{TextEncoder}: FC($1000\!\to\!64$)\,$\to$\,LayerNorm\,$\to$\,
GELU\,$\to$\,Dropout(0.6)\,$\to$\,FC($64\!\to\!64$).

\subsection{Cross-Modal Attention}

Let $\mathbf{h}_m\in\mathbb{R}^d$ denote the encoded representation
of modality $m\in\{f,a,t\}$ (face, audio, text). For each modality,
we compute attention over the other two:

\begin{equation}
\boldsymbol{\alpha}_m =
  \operatorname{softmax}\!\!\left(\frac{\mathbf{W}_Q\mathbf{h}_m\cdot
  \mathbf{K}_{-m}^\top}{\sqrt{d}}\right),
\quad
\mathbf{c}_m = \boldsymbol{\alpha}_m \mathbf{V}_{-m}
\label{eq:attention}
\end{equation}

where $\mathbf{K}_{-m},\mathbf{V}_{-m}\in\mathbb{R}^{(M-1)\times d}$
stack the key and value projections from all other modalities
($M=3$). A residual connection and LayerNorm stabilize training:

\begin{equation}
\tilde{\mathbf{h}}_m = \operatorname{LayerNorm}(\mathbf{c}_m + \mathbf{h}_m)
\label{eq:residual}
\end{equation}

The fused representation is the concatenation:
$\mathbf{z} = [\tilde{\mathbf{h}}_f \| \tilde{\mathbf{h}}_a \|
\tilde{\mathbf{h}}_t] \in \mathbb{R}^{3d=192}$.

\subsection{Multi-Task Prediction Heads}

A shared layer $\mathbf{s}=\operatorname{FC}(192\!\to\!192)
\to\operatorname{LayerNorm}\to\operatorname{GELU}\to\operatorname{Dropout}(0.6)$
feeds three independent heads:

\textbf{Binary head}: $\operatorname{FC}(192\!\to\!2)$ $\to$ logits
$\to$ softmax $P(\text{depressed})$.

\textbf{Score head}: $\operatorname{FC}(192\!\to\!64)
\to\operatorname{GELU}\to\operatorname{Dropout}(0.3)
\to\operatorname{FC}(64\!\to\!1)\to\sigma$,
predicts normalized PHQ-8 score $\hat{s}\in[0,1]$, scaled to $[0,24]$.

\textbf{Symptom head}: $\operatorname{FC}(192\!\to\!128)
\to\operatorname{GELU}\to\operatorname{Dropout}(0.3)
\to\operatorname{FC}(128\!\to\!32)$, reshaped to $(B,8,4)$
for eight 4-class symptom ordinal predictions.

\subsection{Modality Dropout}

During training, each modality is independently zeroed with
probability $p=0.15$, with one Bernoulli draw per modality per
mini-batch (i.e., the drop/keep decision is shared across all
examples in a batch, not sampled independently per example):

\begin{equation}
\mathbf{h}_m^\prime =
  \mathbf{h}_m \cdot \mathbf{1}[\operatorname{Bernoulli}(1-p)=1]_{\text{per batch}}
\label{eq:moddrop}
\end{equation}

This is intended to encourage cross-modal attention to function when
one channel is absent---a real clinical deployment constraint when
video capture or audio recording fails---though because masking is
batch-level rather than per-example, it is a weaker and coarser
regularizer than fully independent per-example modality dropout would
be; this is consistent with, and may partly explain, the mixed (not
uniformly graceful) test-time robustness result in
Section~\ref{sec:robustness}. Total parameters: 592,227.

% ============================================================
\section{Training Objective}
\label{sec:training}

\subsection{Multi-Task Loss}

\begin{equation}
\mathcal{L}_\text{MT} =
  \underbrace{\mathcal{L}_\text{CE}(\hat{y},y;\mathbf{w})}_{\text{binary}}
  + \lambda_s\!\underbrace{\mathcal{L}_\text{MSE}(\hat{s},s/24)}_{\text{PHQ-8 score}}
  + \lambda_y\!\underbrace{\frac{1}{8}\sum_{k=1}^{8}
    \mathcal{L}_\text{CE}(\hat{y}_k,y_k)}_{\text{8 symptoms}}
\label{eq:mt_loss}
\end{equation}

where $\mathbf{w}$ are class weights (capped at $1.8\times$),
$\lambda_s=0.3$, $\lambda_y=0.2$. Auxiliary terms are linearly
warmed up from 0 over epochs~1--15 to ensure the binary head
stabilizes before multi-task gradients become active. Loss weights
($\lambda_s$, $\lambda_y$, $\lambda_f$), dropout rates, and warmup
schedules were selected by a small manual search over
$\{0.1,0.2,0.3,0.5\}$ (auxiliary weights) and
$\{0.3,0.4,0.5,0.6\}$ (dropout) on the dev split only, guided by
dev macro-F1; the train and test splits were never used for this
selection. We report the single configuration selected this way
rather than a full grid, and note that dev-set-guided selection
on $N{=}34$ carries a risk of mild optimistic bias in the reported
dev metrics---a limitation shared with the vast majority of
small-clinical-dataset studies and mitigated here by the
held-out test-split evaluation (Table~\ref{tab:sota}).

\subsection{Equalized Odds Fairness Loss}

\begin{equation}
\mathcal{L}_\text{fair} =
  \underbrace{(\operatorname{TPR}_M - \operatorname{TPR}_F)^2}_{\text{equal sensitivity}}
  + \underbrace{(\operatorname{FPR}_M - \operatorname{FPR}_F)^2}_{\text{equal specificity}}
\label{eq:fairness}
\end{equation}

\begin{equation}
\mathcal{L}_\text{total} = \mathcal{L}_\text{MT} +
  \lambda_f \mathcal{L}_\text{fair}, \quad \lambda_f=0.1
\label{eq:total}
\end{equation}

$\lambda_f$ is activated after 8 warmup epochs to ensure classifier
stability before the fairness gradient becomes active. The FPR term
in~\eqref{eq:fairness} is critical: a model predicting positive for
all participants satisfies TPR equality trivially while maximizing FPR;
the FPR penalty directly prevents this collapse mode.

\subsection{Decision Threshold}
\label{subsec:threshold}

The single-task \texttt{ConfigurableFusionModel} (used for every row
of Table~\ref{tab:ablation} except the proposed model, and for the
scaled significance tests) and the proposed
\texttt{MultiTaskFusionModel} use two different, both dev-safe,
threshold-selection procedures, which we state explicitly here to
avoid ambiguity: the single-task model's threshold is tuned directly
on dev probabilities after loading the best checkpoint (consistent
with Section~\ref{sec:dataset}'s statement that threshold tuning is
performed on the dev split); the proposed multi-task model instead
tunes its threshold on \emph{training}-set probabilities after
loading the best checkpoint,

\begin{equation}
\tau^* = \arg\max_{\tau \in [0.2, 0.8]}
  F_1\!\left(\mathbf{1}[\hat{p} \geq \tau],\, y_{\text{train}}\right)
\label{eq:threshold}
\end{equation}

and applies that threshold fixed to the dev set without re-tuning.
Both procedures use only labels the model is licensed to see before
the held-out test split is touched---the training-set variant simply
substitutes train labels for dev labels as the tuning target---so
neither introduces dev- or test-label leakage, but they are not the
same procedure, and Table~\ref{tab:ablation}'s single-task and
multi-task rows should not be assumed to use identical
threshold-selection logic. We use the training-set variant for the
multi-task model specifically to avoid a threshold–checkpoint
mismatch that early experiments showed could occur if the dev
threshold were taken from the final epoch rather than the best-F1
epoch. A further clarification applies to the downstream analyses in
Sections~\ref{sec:robustness} and~\ref{sec:calibration}: those
experiments each independently re-derive a threshold on dev
probabilities specific to that analysis (e.g., a separate threshold
per masking condition in Section~\ref{sec:robustness}), rather than
reusing the single training-set-tuned threshold above. Every
procedure in this paper selects its threshold using only train and/or
dev labels and applies it fixed to test, so none leaks test labels,
but "the decision threshold" is not one single constant reused
identically everywhere in the paper, and we name which procedure
applies at each point to avoid that impression.

\subsection{Optimization Details}

AdamW: lr\,=\,$5\!\times\!10^{-4}$, weight\_decay\,=\,$10^{-2}$.
CosineAnnealingLR scheduler. Batch size~16. Early stopping:
patience~15, monitoring dev macro-F1 (clinical metric of interest).
Embedding dimension $d=64$ was chosen to prevent memorization on
$N=107$: at $d=256$ (7.8M~params), the model memorized the training
set within 5~epochs (dev AUC~=~0.49). Reducing to $d=64$ (592K~params)
with Dropout(0.6) and weight\_decay\,=\,$10^{-2}$ restored AUC~=~0.66.

% ============================================================
\section{Ablation Study}
\label{sec:ablation}

\subsection{Controlled Modality Ablation}

To isolate modality contributions from architectural confounds, we
train a single \texttt{ConfigurableFusionModel} class across all seven
modality subsets, preserving identical encoders, attention mechanism,
classifier, and optimizer. Each configuration is trained over five
random seeds; we report mean$\pm$std AUC and macro-F1, bootstrap
$\text{CI}_{95\%}$ (2,000~resamples), and paired bootstrap significance.
The five random seeds used throughout this ablation are
$\{42, 1, 7, 123, 2024\}$. Table~\ref{tab:ablation} and
Fig.~\ref{fig:ablation} summarize results.

\begin{table}[t]
\centering
\renewcommand{\arraystretch}{1.15}
\caption{Ablation Study: DAIC-WOZ Dev ($N{=}34$), 5~seeds per configuration}
\label{tab:ablation}
\scriptsize
\setlength{\tabcolsep}{1.5pt}
\begin{tabular}{lcccc}
\toprule
\textbf{Configuration} & \textbf{AUC} & \textbf{95\% CI} & \textbf{F1} & \textbf{Acc} \\
\midrule
Face only$^{\ddagger}$ & $0.513\pm0.113$ & $[0.245,0.691]$ & $0.525\pm0.008$ & 0.382 \\
Text only         & $0.561\pm0.043$ & $[0.379,0.789]$ & $0.539\pm0.018$ & 0.459 \\
Face + Text       & $0.667\pm0.069$ & $[0.636,0.945]$ & $0.570\pm0.063$ & 0.606 \\
Audio + Text      & $0.664\pm0.057$ & $[0.417,0.831]$ & $0.437\pm0.114$ & 0.588 \\
Face + Audio      & $0.712\pm0.028$ & $[0.480,0.883]$ & $0.563\pm0.041$ & 0.612 \\
Audio only        & $\mathbf{0.721\pm0.013}$ & $[0.541,0.899]$ & $0.593\pm0.028$ & 0.676 \\
Tri-Modal (ST)    & $0.699\pm0.047$ & $[0.449,0.863]$ & $0.607\pm0.063$ & 0.688 \\
\midrule
\textbf{Proposed (MT)} & $0.658\pm0.029$ & $[0.500,0.879]$
  & $\mathbf{0.629\pm0.021}$ & $\mathbf{0.688}$ \\
\bottomrule
\multicolumn{5}{p{0.95\columnwidth}}{\scriptsize ST~=~single-task,
MT~=~multi-task. All rows except Proposed~(MT) use 120-dim
MFCC-only audio; the Proposed~(MT) row uses the extended 199-dim
MFCC+COVAREP+FORMANT audio (Section~\ref{sec:features}).
$^{\ddagger}$The
face-only row reflects the feature pipeline used throughout this table.
Section~\ref{sec:sentinel} shows this configuration is depressed to chance
by an OpenFace $-100$ missing-data sentinel; with sentinel-carrying frames
removed it reaches AUC~$0.708\pm0.042$. The remaining rows are unaffected
in practice---re-running the tri-modal configuration with the corrected
face features changes AUC by $+0.005$ ($p{=}0.313$)---so this table is
reported as-run for consistency with the rest of the paper.} \\
\end{tabular}
\end{table}

\begin{figure}[t]
  \centering
  \includegraphics[width=\columnwidth]{fig2_ablation_ci.png}
  \caption{Per-configuration AUC (mean~$\pm$~std, 5~seeds). Tri-modal
    achieves the best macro-F1~$=0.607$; the proposed multi-task model
    further improves to F1~$=0.629\pm0.021$.}
  \label{fig:ablation}
\end{figure}

\noindent\textbf{Statistical note:} The proposed multi-task model
achieves mean F1~$=0.629\pm0.021$ at the originally-reported 5-seed
sample, a $+0.022$ improvement in mean over the single-task tri-modal
baseline ($0.607\pm0.063$), with markedly reduced seed variance
($\pm0.021$ vs.\ $\pm0.063$), and all five multi-task seeds
individually exceed the single-task baseline's \emph{mean} of $0.607$
(Fig.~\ref{fig:seeds}). This is a weaker claim than paired
seed-for-seed superiority, so we additionally recovered the
previously-unreported per-seed single-task F1 values and ran a paired
Wilcoxon signed-rank test matching each multi-task run against the
single-task run from the \emph{same} seed. At $n{=}5$ that test did
\emph{not} reach significance ($p{=}0.625$): the multi-task model won
only 3 of the 5 matched pairs. With only $n{=}5$ paired seeds, an
exact two-sided Wilcoxon signed-rank test cannot reach $p<0.05$ under
\emph{any} possible outcome, so this result was structurally
uninformative rather than a genuine null. A power analysis of the
5-seed differences (Cohen's~$d{=}0.34$) indicated $n{\approx}67$
paired seeds would be needed for 80\% power at $\alpha{=}0.05$; we
therefore scaled the comparison to $n{=}70$ paired seeds (seeds
$1$--$70$), re-training both arms from scratch at each seed.

The scaled test resolves the question decisively. Across 70 seeds,
the multi-task model remains stable and strong (mean
F1~$=0.657\pm0.037$), while the single-task baseline shows somewhat
higher variance than the original 5-seed sample implied (mean
F1~$=0.597\pm0.076$, range $[0.46,0.77]$ across seeds, vs.\ the
originally-reported $0.607\pm0.063$): several single-task seeds
show reduced F1 despite acceptable AUC ($0.60$--$0.72$), suggesting
the single-task model's discriminative ranking is often fine but its
decision threshold is more seed-sensitive on this small dev set than
the multi-task heads and modality dropout, which regularize this
partially away. The paired Wilcoxon test is now unambiguous: multi-task
wins 51 of 70 seed pairs (Wilcoxon $p{=}8.1{\times}10^{-7}$,
Cohen's~$d{=}0.72$)---a medium-to-large effect that is in fact larger
than the $d{=}0.34$ assumed for the original power calculation,
meaning that calculation was itself conservative and fewer than the
estimated 67~seeds would likely have sufficed; scaling to 70 removes
any doubt. The original 5-seed mean ($0.607\pm0.063$) is close to,
not dramatically below, the 70-seed mean, so its main limitation was
that $n{=}5$ is structurally too few for an exact Wilcoxon test to
reach significance (Section~\ref{sec:ablation} above), not that the
mean itself was badly biased. This is nonetheless a concrete,
within-pipeline illustration of the seed-sampling risk this paper
otherwise raises as an external methodological concern
\cite{bouthillier2021}: small seed counts can leave a real,
confirmable effect looking merely suggestive.

Taken together, these results show multi-task learning is a
statistically confirmed result at this sample size, not merely an
encouraging trend. We note this comparison still bundles
two simultaneous changes---the multi-task prediction heads and the
extended 199-dim MFCC+COVAREP+FORMANT audio features, versus the
single-task baseline's 120-dim MFCC-only audio---so this result
confirms the combined intervention rather than isolating multi-task
learning alone as the causal factor; disentangling the two would
require a further controlled ablation, which we leave to future work.
The cross-modal-attention and fairness-loss comparisons in
Section~\ref{sec:fairness} remain at $n{=}10$ and are not rescaled
here, so they retain the inconclusive-rather-than-null status
described above. Fig.~\ref{fig:p7p8} summarizes both the original
5-seed sample and the scaled 70-seed test side by side, and
Fig.~\ref{fig:seeds} shows the full 70-seed paired distribution.

\noindent\textbf{Multiple-comparisons correction:} The three design
choices evaluated in this testbed---multi-task learning, cross-modal
attention, and the Equalized Odds fairness loss---are naturally
treated as one family of comparisons, since all three test whether a
proposed architectural addition improves on a shared tri-modal
baseline. Applying a Bonferroni correction across this family of
three ($\alpha=0.05/3=0.0167$) does not change any conclusion: the
scaled multi-task result remains significant by a wide margin
($p{=}8.1{\times}10^{-7}\times3=2.4{\times}10^{-6}$, still far below
$\alpha$), while the attention ($p{=}0.496$) and fairness ($p{=}0.813$)
comparisons remain non-significant under either the corrected or
uncorrected threshold. We report this correction for completeness
rather than because it changes the substantive finding: at the
sample sizes and effect sizes involved here, multi-task learning
would survive a far stricter correction than Bonferroni, and the two
inconclusive comparisons were already reported as such.

\begin{figure}[t]
  \centering
  \includegraphics[width=\columnwidth]{fig4_seed_consistency.png}
  \caption{Single-task baseline vs.\ proposed multi-task model F1,
    box plot with individual seeds overlaid, across the full scaled
    test ($n{=}70$ paired seeds). The single-task baseline shows
    somewhat higher seed-to-seed spread (range $[0.46,0.77]$, mean
    $0.597\pm0.076$) than the multi-task model's tighter cluster
    (mean $0.657\pm0.037$), consistent with the Statistical note. At
    the original $n{=}5$ sample this comparison was inconclusive
    ($p{=}0.625$); the full $n{=}70$ test is decisive
    ($p{=}8.1{\times}10^{-7}$, $d{=}0.72$, wins 51/70).}
  \label{fig:seeds}
\end{figure}

\begin{figure}[t]
  \centering
  \includegraphics[width=\columnwidth]{fig3_p7_vs_p8.png}
  \caption{Single-task baseline vs.\ proposed multi-task model F1,
    mean\,$\pm$\,std. \emph{Left:} the original 5-seed sample
    ($0.607\pm0.063\to0.629\pm0.021$, underpowered paired test,
    $p{=}0.625$). \emph{Right:} the scaled 70-seed test
    ($0.597\pm0.076\to0.657\pm0.037$, confirmed,
    $p{=}8.1{\times}10^{-7}$). The 70-seed single-task mean is close
    to the original 5-seed estimate; the decisive difference is that
    $n{=}70$ gives the exact Wilcoxon test enough seeds to reach
    significance, which $n{=}5$ cannot under any outcome.}
  \label{fig:p7p8}
\end{figure}

\subsection{Cross-Modal Attention vs.\ Plain Concatenation}
\label{subsec:attn_vs_concat}

A critical reviewer concern is whether cross-modal attention provides
measurable benefit over simpler feature concatenation. To answer this
with controlled evidence, we implement \texttt{ConcatFusionModel}:
identical encoders, identical classifier, but the fusion step replaces
Eq.~\eqref{eq:attention}--\eqref{eq:residual} with
$\mathbf{z}=[\mathbf{h}_f\|\mathbf{h}_a\|\mathbf{h}_t]$.
Both systems are trained under identical hyperparameters over
\textbf{10 random seeds} $\{42,1,7,13,21,100,2024,7777,555,88\}$
(Table~\ref{tab:attn_concat}, Fig.~\ref{fig:attn_concat}).

\begin{table}[t]
\centering
\renewcommand{\arraystretch}{1.15}
\caption{Cross-Modal Attention vs.\ Plain Concatenation (10~seeds, paired Wilcoxon test)}
\label{tab:attn_concat}
\scriptsize
\setlength{\tabcolsep}{0.5pt}
\begin{tabular}{lcccc}
\toprule
\textbf{Fusion} & \textbf{Mean F1} & \textbf{Mean AUC}
  & \textbf{F1 wins} & \textbf{AUC wins} \\
\midrule
Cross-Modal Attention & $0.668\pm0.050$ & $0.667\pm0.064$
  & 4/10 & \textbf{6/10} \\
Plain Concatenation   & $0.644\pm0.033$ & $0.655\pm0.046$
  & \textbf{5/10} & 4/10 \\
\midrule
$\Delta$ (Attn $-$ Concat) & $+0.024$ & $+0.012$ & & \\
Wilcoxon $p$ & 0.496 & 0.647 & \multicolumn{2}{c}{n.s.\ at $\alpha{=}0.05$} \\
\bottomrule
\multicolumn{5}{l}{\scriptsize 1 tie in F1. n.s. = not significant.} \\
\end{tabular}
\end{table}

\begin{figure}[t]
  \centering
  \includegraphics[width=\columnwidth]{fig8_attention_vs_concat.png}
  \caption{Cross-modal attention vs.\ plain concatenation (10~seeds,
    identical encoders/classifier). Attention achieves higher mean F1
    (+0.024) and higher mean AUC (+0.012), winning 6/10~seeds on AUC,
    but neither difference is statistically significant
    (Wilcoxon $p{=}0.496$/$0.647$, $N{=}34$).}
  \label{fig:attn_concat}
\end{figure}

Cross-modal attention achieves mean F1\,$=0.668\pm0.050$ vs.\
concat $0.644\pm0.033$ ($\Delta=+0.024$) and mean AUC\,$=0.667$
vs.\ $0.655$ ($\Delta=+0.012$). Attention wins 6/10~seeds on AUC.
However, neither difference reaches statistical significance
(Wilcoxon signed-rank $p{=}0.496$ for F1; $p{=}0.647$ for AUC, $n{=}10$
paired seeds), due to high seed variance on the small dev set
($N{=}34$). To characterize \emph{why} the test is inconclusive rather
than simply reporting non-significance, we compute the paired effect
size (Cohen's~$d$) and a post-hoc power analysis: F1 shows a small-to-
medium effect ($d{=}0.34$), requiring $n{\approx}69$ paired seeds for
80\% power at $\alpha{=}0.05$ (current power at $n{=}10$: 16\%); AUC
shows a small effect ($d{=}0.19$, $n_{80\%}{\approx}211$, current
power~9\%). This confirms the study is genuinely underpowered to
detect an effect of this size, rather than providing evidence of no
effect. We retain cross-modal attention as the default fusion
mechanism for two reasons: (1)~a consistent positive trend in both
metrics across 10~seeds, and (2)~the interpretability benefit of
attention weights for clinical auditing, which concatenation does not
provide. Definitive superiority claims require a substantially larger
paired-seed study or a larger dataset.

\subsection{Pooled 10-Fold Cross-Validation}
\label{sec:kfold}

The headline metrics above come from the single fixed AVEC-2017
split (train=107 / dev=34 / test=47). A single split on $N{=}34$--$47$
has limited statistical reliability, a general concern for small
clinical cohorts. To obtain a larger effective
held-out sample, we pool all 188 labeled DAIC-WOZ participants
(train+dev+test) and run stratified $k$-fold cross-validation with
$k{=}5$ under two independent partition seeds ($\{42,7\}$), giving
10~fold-runs in total, each with an internal 15\% validation slice
carved from the training portion for early stopping and threshold
selection (identical training protocol to
Section~\ref{sec:training}). Test-origin participants contribute
only binary label, PHQ-8 total score, and gender (no per-symptom
items in the public release), so the symptom-classification loss
term is masked for those samples during cross-validation; all 188
participants contribute to the binary and PHQ-8-score losses.

\begin{table}[t]
\centering
\renewcommand{\arraystretch}{1.15}
\caption{Pooled 10-Fold Cross-Validation ($N{=}188$, 2~partition seeds
$\times$ 5~folds)}
\label{tab:kfold}
\footnotesize
\begin{tabular}{lc}
\toprule
\textbf{Metric} & \textbf{Mean $\pm$ Std (across 10 folds)} \\
\midrule
Macro-F1  & $0.554 \pm 0.093$ \\
AUC-ROC   & $0.612 \pm 0.120$ \\
Accuracy  & $0.657 \pm 0.058$ \\
TPR gap (male$-$female) & $0.415 \pm 0.211$ \\
\bottomrule
\end{tabular}
\end{table}

The cross-validated estimates are lower than the fixed-split dev/test
numbers reported above (F1 $0.554$ vs.\ $0.629$/$0.585$; AUC $0.612$
vs.\ $0.658$/$0.706$), and the fold-to-fold spread is substantial
(F1 std $0.093$, AUC std $0.120$). Most notably, the fairness picture
changes: the mean TPR gap across folds is $0.415\pm0.211$, and in
7~of the 10~folds the model achieves $\text{TPR}_\text{female}=0$
(no depressed female participant in that fold's test partition is
correctly flagged), compared to the single-split test result of
TPR~gap\,$=0.000$ reported in Table~\ref{tab:sota}. We interpret this
as evidence that the fixed-split fairness result is not a stable
property of the model but an artifact of one particular
train/dev/test partition; the cross-validated estimate is the more
trustworthy characterization of both overall performance and fairness
behavior on this dataset at this sample size. We report both sets of
numbers---fixed-split and cross-validated---rather than presenting
only the more favorable one.

% ============================================================
\section{Comparison with Published Methods}
\label{sec:sota}

\begin{table}[t]
\centering
\renewcommand{\arraystretch}{1.15}
\caption{Comparison with Published DAIC-WOZ Results}
\label{tab:sota}
\scriptsize
\setlength{\tabcolsep}{1.5pt}
\begin{tabular}{lccccc}
\toprule
\textbf{Method} & \textbf{Mod.} & \textbf{Split} & \textbf{F1} & \textbf{AUC} & \textbf{TPR Gap} \\
\midrule
AVEC-2017 audio \cite{ringeval2017}      & A     & Dev & 0.50 & --- & --- \\
AVEC-2017 text \cite{ringeval2017}       & T     & Dev & 0.49 & --- & --- \\
Williamson et al.\ \cite{williamson2016} & A+V   & Dev & 0.57 & --- & --- \\
Gong \& Poellabauer \cite{gong2017}      & A+V+T & Dev & 0.70 & --- & --- \\
Ray et al.\ \cite{ray2019}               & A+T   & Dev & --- & 0.80 & --- \\
\textbf{This work} & \textbf{F+A+T} & \textbf{Dev}
  & $\mathbf{0.629}$ & $0.658$ & 0.286$^{\ddagger}$ \\
\textbf{This work} & \textbf{F+A+T} & \textbf{Test}
  & $\mathbf{0.585}$ & $\mathbf{0.706}$ & $\mathbf{0.000}^{\dagger}$ \\
\bottomrule
\multicolumn{6}{p{0.95\columnwidth}}{\scriptsize Mod.=Modalities;
A=Audio, V=Video, T=Text, F=Face AUs. `---' = not reported by
authors. Dev row: F1~$0.629\pm0.021$, AUC~$0.658\pm0.029$; Test row
($N{=}47$).} \\
\multicolumn{6}{p{0.95\columnwidth}}{\scriptsize $^{\ddagger}$The Dev
F1/AUC use the headline train-tuned threshold
(Section~\ref{subsec:threshold}); the Dev TPR Gap uses the fairness
audit's independently dev-tuned threshold ($\tau{=}0.325$,
Table~\ref{tab:fairness}), so these two Dev-row entries are not from
a single shared threshold.} \\
\multicolumn{6}{p{0.95\columnwidth}}{\scriptsize $^{\dagger}$Equal but low TPR (0.286 both groups); 10-fold CV gives $0.415\pm0.211$ (Sec.~\ref{sec:kfold}).} \\
\end{tabular}
\end{table}

\begin{figure}[t]
  \centering
  \includegraphics[width=\columnwidth]{fig5_sota.png}
  \caption{DAIC-WOZ dev F1: this work vs.\ classical (non-transformer)
    published baselines. Ray et al.\ \cite{ray2019} report higher AUC on different metrics/splits
    and are not directly comparable on this figure (see Table~\ref{tab:sota}).}
  \label{fig:sota}
\end{figure}

Our system achieves dev F1\,$=0.629\pm0.021$, outperforming three older
baselines (AVEC-2017 audio, AVEC-2017 text, Williamson et al.) on F1,
while remaining below Gong \& Poellabauer's topic-modeling result
(F1\,$=0.70$); see Table~\ref{tab:sota} for the full comparison.
More recent transformer-based methods \cite{ray2019} achieve
higher AUC ($\approx$0.80) using pre-trained language models and larger
GPU-trained architectures; our system does not claim to surpass these.
Our contribution is methodological: controlled ablation with an
identical model class, symptom-level multi-task learning, and honest
characterization of failure modes (fairness loss, attention), rather
than raw performance advancement. Fig.~\ref{fig:sota} visualizes this
positioning against the classical (non-transformer) baselines.

On the held-out test split ($N{=}47$), the model achieves
F1\,$=0.585$ and AUC\,$=0.706$---the latter \emph{exceeding} the dev
AUC of $0.658\pm0.029$, indicating the model generalizes rather than
overfitting to the dev set. The test F1 drop ($0.629\to0.585$) is
expected: per Section~\ref{subsec:threshold}, the multi-task model's
decision threshold is tuned on the training set and may not
optimally calibrate for the test distribution. The test-set gender
fairness audit reveals a TPR gap of~0.000: male and female depressed
participants are detected at \emph{equal} rates, but that rate is only
$0.286$ for both groups (4/14 depressed test participants correctly
flagged)---the model is missing most depressed participants in this
split \emph{regardless of gender}, so equal TPR here reflects equally
poor sensitivity rather than a fairness success. Pooled 10-fold
cross-validation (Section~\ref{sec:kfold}) confirms this single-split
gap of~0.000 is not representative: across 188 participants the mean
TPR gap is $0.415\pm0.211$, with several folds showing $\text{TPR}_\text{female}=0$
entirely. We therefore do not treat the fixed-split TPR gap as
evidence of fairness and report the cross-validated estimate as the
more reliable characterization.

% ============================================================
\section{Fairness Audit}
\label{sec:fairness}

\subsection{Four-Criterion Gender Audit}

Following Barocas et al.\ \cite{barocas2019}, we evaluate four
group-fairness criteria between male ($M$) and female ($F$)
participants on the dev set (Table~\ref{tab:fairness},
Fig.~\ref{fig:radar}).

\begin{table}[t]
\centering
\renewcommand{\arraystretch}{1.15}
\caption{Gender Fairness Audit (DAIC-WOZ Dev, $N{=}34$)}
\label{tab:fairness}
\scriptsize
\setlength{\tabcolsep}{2pt}
\begin{tabular}{lcccc}
\toprule
\textbf{Criterion} & \textbf{Male} & \textbf{Female} & \textbf{Gap}
  & \textbf{Verdict} \\
\midrule
Equal Opportunity (TPR) & 0.71 & 1.00 & \textbf{0.286}
  & \textcolor{red}{\textbf{Biased}} \\
Predictive Parity (PPV) & 0.45 & 0.45 & 0.000 & Fair \\
Demographic Parity      & 0.61 & 0.69 & 0.076 & Fair \\
F1-score Parity         & 0.56 & 0.63 & 0.069 & Fair \\
\bottomrule
\multicolumn{5}{p{0.95\columnwidth}}{\scriptsize This audit uses a
decision threshold ($\tau{=}0.325$) selected independently for this
analysis, per the threshold-selection note in
Section~\ref{subsec:threshold}; it is not the same threshold as the
$\tau{=}0.48$ used for the headline calibration and error-case
analysis in Section~\ref{sec:calibration}, so per-group FN/TP counts
here (male FN~$=2$, female FN~$=0$) are not directly comparable to
that section's dev-wide FN~$=8$, which reflects a stricter
threshold.} \\
\end{tabular}
\end{table}

\begin{figure}[t]
  \centering
  \includegraphics[width=0.85\columnwidth]{fig6_fairness_radar.png}
  \caption{Gender fairness radar (four criteria). Equal-Opportunity
    gap~=~0.286 is the only criterion exceeding the 0.10 fairness threshold,
    motivating the Equalized Odds training loss.}
  \label{fig:radar}
\end{figure}

The Equal-Opportunity gap of~0.286 reveals that the model correctly
detects depression in 100\% of depressed women but only 71\% of
depressed men. A clinical tool missing 29\% of depressed men is
unacceptable---this gap motivates the Equalized Odds constraint as
a \emph{clinically necessary} design choice.

\subsection{Fairness Loss Effectiveness}

To empirically evaluate $\mathcal{L}_\text{fair}$, we train with
$\lambda_f=0$ (no fairness) versus $\lambda_f=0.1$ over \textbf{10 seeds}
$\{42,1,7,13,21,100,2024,7777,555,88\}$
(Table~\ref{tab:fairness_exp}, Fig.~\ref{fig:fairness_ba}).

\begin{table}[t]
\centering
\renewcommand{\arraystretch}{1.15}
\caption{Equalized Odds Loss: 10-Seed Paired Analysis}
\label{tab:fairness_exp}
\scriptsize
\setlength{\tabcolsep}{0.5pt}
\begin{tabular}{lcccc}
\toprule
\textbf{Condition} & \textbf{TPR Gap} & \textbf{Wins}
  & \textbf{F1} & \textbf{$p$} \\
\midrule
No fairness ($\lambda_f{=}0$)     & $0.320\pm0.191$ & 4/10 & $0.668\pm0.050$ & \multirow{2}{*}{0.813} \\
With fairness ($\lambda_f{=}0.1$) & $0.311\pm0.233$ & 3/10 & $0.666\pm0.038$ & \\
\midrule
$\Delta$ (with $-$ without) & $-0.009$ & \multicolumn{2}{c}{3 ties} & n.s. \\
\bottomrule
\multicolumn{5}{l}{\scriptsize n.s. = not significant at $\alpha{=}0.05$. F1 cost: $-0.001$.} \\
\end{tabular}
\end{table}

\begin{figure}[t]
  \centering
  \includegraphics[width=\columnwidth]{fig7_fairness_before_after.png}
  \caption{Per-seed TPR gap with/without Equalized Odds loss (10~seeds).
    Mean gap: $0.311$ (with) vs.\ $0.320$ (without); Wilcoxon $p{=}0.813$.
    The fairness loss fails to significantly reduce bias at $N{=}34$.}
  \label{fig:fairness_ba}
\end{figure}

\textbf{Negative result:} The 10-seed paired analysis reveals that
$\mathcal{L}_\text{fair}$ does \emph{not} significantly reduce the
TPR gap at this sample size (Wilcoxon $p{=}0.813$). Fairness wins 3/10
seeds, the baseline wins 4/10, with 3~ties. The worst-case TPR gap
\emph{increased} in the high-variance seed
(no-fairness $0.714$ $\to$ with-fairness $0.857$), counter to the
intended effect. The F1 cost is negligible ($-0.001$).

We report this negative result transparently: bias mitigation via a
differentiable surrogate loss is insufficient at $N{=}34$, where
the gender-stratified subgroups contain as few as 5~participants.
The paired effect size is negligible (Cohen's~$d{=}-0.07$), and a
post-hoc power analysis shows that even a study with $n{\approx}1{,}500$
paired seeds would be needed to reach 80\% power for an effect this
small at $\alpha{=}0.05$ (current power at $n{=}10$: 6\%)---indicating
the true effect, if any, is likely too small to be practically
meaningful at this scale, rather than merely undetected. This finding
aligns with the theoretical result of \cite{barocas2019} that fairness
constraints require adequate representation in both groups to converge
reliably. The Equalized Odds loss remains in our framework as a
principled design choice and may become effective with larger clinical
datasets or balanced gender splits---but we make no claim of empirical
effectiveness at the current scale.

% ============================================================
\section{Explainability}
\label{sec:xai}

Three mechanistically independent attribution methods are applied to
the trained model. Convergence across independent methods is stronger
evidence than any single technique alone \cite{abnar2020}.

\begin{figure}[t]
  \centering
  \includegraphics[width=\columnwidth]{fig13_explainability.png}
  \caption{Three convergent explainability methods. (a)~Attention
    rollout ranks modality importance (audio highest, face and text
    close behind). (b)~Gradient$\times$input saliency over the top-6
    attributed AU channels: AU04 (highlighted) ranks fifth, within
    the top~5 but not top-ranked. (c)~Occlusion ablation: removing
    audio or text costs $\approx$4$\times$ the AUC that removing face
    does, consistent with (a).}
  \label{fig:explainability}
\end{figure}

\textbf{Cross-modal attention rollout} \cite{abnar2020}:
Audio receives the highest mean attention weight~(0.39), followed
by text~(0.31) and face~(0.30). This is consistent with the
ablation result that audio is the strongest single modality.

\textbf{Gradient\,$\times$\,input saliency} \cite{shrikumar2017}:
$\|\nabla_{\mathbf{x}} f \odot \mathbf{x}\|_1$ attributed over
20~AU channels ranks \textbf{AU20 (Lip Stretcher)} highest
(importance~$0.130$), followed by AU06 (Cheek Raiser, $0.120$), AU05
(Upper Lid Raiser, $0.120$), AU25 (Lips Part, $0.085$), and
\textbf{AU04 (Brow Lowerer, $0.079$)} fifth. AU04 is a FACS unit
established in the clinical literature as a depressed-affect marker
\cite{ekman1978} and does appear within the top~5 attributed
channels, but it is not the single top-ranked feature by this
method---AU20, the actual top-ranked channel, does not have the same
established link to depressed affect in the FACS literature we
reviewed. We report the full ranking rather than only the clinically
convenient part of it.

Two caveats apply to this ranking. First, all five top-attributed
channels are continuous \emph{intensity} channels ($r$ suffix), which
Section~\ref{sec:sentinel} confirms are free of the $-100$ sentinel
(the sentinel occurs only in the six binary presence channels), so the
ranking is not an artefact of that fault; the attributions are
nonetheless computed from a model trained with the six affected
channels present, and we report them as such. Second, AU04, AU12 and
AU15 each occur twice in the channel set---once as intensity, once as
presence---so attributions must be keyed by channel type rather than by
bare AU name; keying on the name alone silently collapses the pair
(for AU04, the presence-channel value $0.011$ would displace the
intensity-channel value $0.079$ that the ranking above refers to).

\textbf{Leave-one-modality-out AUC drop}: Removing audio causes the
largest AUC decrease ($-0.03$) on the dev split, consistent with
audio dominance via occlusion without relying on gradient computation.
This dev-split finding does not generalize cleanly to the test split,
however: the test-time masking experiment in
Section~\ref{sec:robustness} finds that removing audio does
\emph{not} hurt test AUC (it improves slightly), while removing text
does---so "most-attended/most-occlusion-sensitive on dev" and
"most necessary for test-set generalization" are shown by this paper
to be distinct properties, not interchangeable evidence of a single
modality's importance.

All three methods consistently rank audio as the most informative
modality \emph{on the dev split}, and gradient saliency places the
clinically-established marker AU04 within its top~5 attributed facial
channels, though not as the single top-ranked feature (AU20 ranks
higher and lacks the same established clinical link). We report this
as partial, not clean, convergence with clinical priors---the
attribution evidence is directionally plausible but does not single
out AU04 as the decisive top-ranked feature. Convergence
across independent methods is still stronger evidence than any single
technique alone, and Section~\ref{sec:robustness} further shows the
dev-split modality ranking does not fully carry over to test-time
necessity.

% ============================================================
\section{Additional Validation Experiments}
\label{sec:additional_validation}

\subsection{Symptom-Level Prediction Heterogeneity}
\label{sec:symptom}

A recent large-scale reproducibility study argues that DAIC-WOZ
classifiers may largely learn \emph{disorder-general} distress cues
(flat affect, low energy, negative language) rather than markers
specific to individual depressive symptoms, since most published
systems are evaluated only on the aggregate binary label. Because our
model predicts all eight PHQ-8 symptom items individually, we can test
this directly: if the model has learned only one generic distress
signal, presence/absence accuracy should be roughly uniform across
all eight items; if it captures symptom-specific structure, accuracy
should vary meaningfully by symptom.

On the dev split, presence/absence macro-F1 ranges from $0.209$
(Psychomotor Change) to $0.424$ (Tired/Low Energy) across the eight
symptoms---a spread of $0.215$, not the uniform profile expected from
a single generic-distress signal. However, the result is not
unambiguously favorable: the two best-predicted symptoms (Tired/Low
Energy, $0.424$; Feeling Down/Depressed, $0.393$) are themselves the
most generic affective/energy markers, while more specific items
(Psychomotor Change, $0.209$; Concentration Problems and No
Interest/Pleasure, both $0.320$) are predicted worse than average. We
therefore report this as partial evidence against a purely
disorder-general account---performance is not uniform---while noting
that the model still performs best on the most generic symptoms,
which is directionally consistent with, not a clean refutation of,
the disorder-general-cues concern. Exact 4-class severity accuracy is
low across all symptoms ($0.029$--$0.412$), so these findings should
be read as evidence about \emph{which symptoms carry a detectable
presence/absence signal}, not as evidence of clinically usable
symptom-severity grading. Section~\ref{sec:discussion}'s
diagnostic-specificity probe tests the same concern from a
complementary, external angle---asking whether the model's depression
score separates depression from other distress conditions in an
independent corpus, rather than asking whether its internal
symptom-level predictions are uniform within DAIC-WOZ---and reaches
the same qualified conclusion: partial, not clean, evidence against a
purely disorder-general account.

\subsection{Missing-Modality Robustness (Test-Time Masking)}
\label{sec:robustness}

Modality dropout is used during training as a regularizer, motivated
by the clinical constraint that a modality may be unavailable at
inference time. To validate this claim directly rather than merely
asserting it, we zero-mask each modality individually at test time
(all-modality-present as baseline) and re-evaluate on both dev and
test splits with the same trained model. Following the paper's stated
protocol, the decision threshold for every masking condition is
selected on the \emph{dev} split's probabilities only and then applied
fixed to the corresponding test-split evaluation---no threshold is
ever selected using test labels. With this protocol, the unmasked
(``None'') baseline test row below reproduces the paper's headline
test F1 exactly ($0.585$), confirming the two evaluations are
consistent.

\begin{table}[t]
\centering
\renewcommand{\arraystretch}{1.15}
\caption{Test-Time Missing-Modality Robustness (threshold selected on
dev only, applied fixed to test)}
\label{tab:robustness}
\begin{tabular}{lcccc}
\toprule
\textbf{Masked} & \multicolumn{2}{c}{\textbf{Dev ($N{=}34$)}}
  & \multicolumn{2}{c}{\textbf{Test ($N{=}47$)}} \\
\textbf{Modality} & \textbf{F1} & \textbf{AUC} & \textbf{F1} & \textbf{AUC} \\
\midrule
None (baseline)  & 0.622 & 0.663 & 0.585 & 0.706 \\
Face masked      & 0.622 & 0.655 & 0.585 & 0.712 \\
Audio masked     & 0.583 & 0.580 & 0.527 & 0.725 \\
Text masked      & 0.555 & 0.610 & 0.415 & \textbf{0.435} \\
\bottomrule
\end{tabular}
\end{table}

The result is a partial, modality-dependent robustness profile rather
than a uniform one. Removing face costs almost nothing on either
split---at the dev-derived threshold, it does not change a single
test-split prediction, consistent with face being the weakest single
modality in the ablation study. Removing audio costs moderately on
both splits by F1 ($0.622\to0.583$ dev, $0.585\to0.527$ test), while
its effect on AUC is mixed (drops on dev, a small, likely-noise
increase on test at $N{=}47$). Removing text is the clearly critical
case on both metrics and both splits: test F1 falls from~$0.585$
to~$0.415$ and test AUC collapses from~$0.706$ to~$0.435$---below
chance---indicating the model's test-set generalization depends
heavily on text being present, even though attention-rollout
(Section~\ref{sec:xai}) ranks audio as the most-attended modality on
average \emph{on dev}. We interpret this as evidence that average
dev-split attention weight and test-time necessity are distinct
properties: a modality can receive lower average attention while
still being critical for cross-modal integration at test time.
Practically, this means the modality-dropout regularizer does
\emph{not} deliver uniform graceful degradation, and a deployed
version of this system should not be assumed safe to run with text
unavailable without further validation.

\subsection{Calibration and Error Analysis}
\label{sec:calibration}

Beyond discrimination (AUC) and thresholded classification (F1), we
assess whether the model's predicted probabilities are meaningful as
probabilities. We compute the Brier score and Expected Calibration
Error (ECE, 5 equal-width bins) on both splits. Dev: Brier
$=0.233$, ECE$=0.102$. Test: Brier$=0.227$, ECE$=0.157$. These
indicate moderate miscalibration---predicted probabilities should not
be read as precise clinical risk estimates without post-hoc
recalibration (e.g., Platt scaling or isotonic regression) on a larger
held-out set.

Error case analysis (dev threshold selected on dev via macro-F1, then
applied fixed to test---no test-label leakage) shows a consistent
false-negative-skewed error profile on \emph{both} splits: on dev
(10/34 misclassified), 8~FN vs.\ 2~FP; on test (14/47 misclassified,
matching the headline test accuracy of $70.2\%$), 10~FN vs.\ 4~FP.
This asymmetry---missing depressed participants more often than
falsely flagging healthy ones---is clinically relevant, since it is
the less safe failure mode for a screening tool, and it holds on both
the tuned split and the independently held-out split rather than
being a dev-specific artifact. This should be weighed when selecting
a deployment threshold (a lower threshold trades false negatives for
more false positives).

To quantify this trade-off concretely, we scan the decision threshold
on the dev split alone (no test-label access) and locate
$\tau{=}0.45$, the smallest threshold before dev macro-F1 begins to
fall sharply. Applied fixed to the test set, exactly as the headline
$\tau{=}0.48$ was, this recall-prioritized operating point raises
test recall on the depressed class from $0.286$ to $0.786$
($4/14\to11/14$ depressed participants correctly flagged) while
macro-F1 is essentially unchanged ($0.585\to0.589$); precision falls
from $0.500$ to $0.407$ (4~FP$\to$16~FP). We do not adopt this as the
headline operating point---it was not selected before this
post-hoc analysis and macro-F1 alone cannot certify it as
generalizing better than $\tau{=}0.48$---but it demonstrates that a
substantially safer operating point (roughly $3\times$ fewer missed
depressed participants) is available near-costlessly in aggregate
F1 terms, simply by moving the threshold, with no retraining. For a
screening deployment where a missed case is more costly than a false
alarm, $\tau{\approx}0.45$ is the more defensible choice; we report
$\tau{=}0.48$ throughout the rest of the paper only for consistency
with the threshold-selection procedure used for the headline result
(Section~\ref{sec:training}), decided before this post-hoc
recall-oriented analysis was run, not because it was formally
pre-registered.

\subsection{Internal Audio Feature Ablation}
\label{sec:audio_ablation}

The clinical audio representation combines three feature families
(MFCC, COVAREP, FORMANT; Section~\ref{sec:features}), but the paper
never isolated which family drives the audio modality's strong
ablation performance. We re-train the single-task tri-modal model
with the audio branch restricted to each family individually and to
pairwise combinations, holding face, text, architecture, and training
protocol fixed (seeds $\{42,1,7\}$, a subset of the 5-seed ablation
protocol, chosen to bound the computational cost of the six additional
configurations; a smaller seed count than elsewhere in this paper, so
these results should be read as suggestive rather than as precisely
estimated as the 5-seed and 10-seed experiments above). This
experiment independently re-trains the full-audio (MFCC+COVAREP+FORMANT)
configuration under the 3-seed protocol as an internal check; its
AUC ($0.711\pm0.055$) is close to but not identical to the 5-seed
Table~\ref{tab:ablation} audio-only figure ($0.721\pm0.013$), the
difference reflecting seed-count variance between the two independent
runs of the same nominal configuration rather than an error.

\begin{table}[t]
\centering
\renewcommand{\arraystretch}{1.15}
\caption{Internal Audio Feature Ablation (Dev, 3~seeds)}
\label{tab:audio_ablation}
\scriptsize
\begin{tabular}{lcc}
\toprule
\textbf{Audio Subset} & \textbf{Dim} & \textbf{AUC (mean$\pm$std)} \\
\midrule
MFCC only              & 120 & $0.710\pm0.053$ \\
COVAREP only           &  74 & $0.619\pm0.084$ \\
FORMANT only           &   5 & $0.688\pm0.060$ \\
MFCC+COVAREP           & 194 & $0.682\pm0.033$ \\
COVAREP+FORMANT        &  79 & $0.569\pm0.039$ \\
MFCC+COVAREP+FORMANT   & 199 & $0.711\pm0.055$ \\
\bottomrule
\end{tabular}
\end{table}

MFCC alone (AUC~$0.710\pm0.053$) achieves essentially the same
discrimination as the full 199-dim combination
(AUC~$0.711\pm0.055$)---the addition of COVAREP and FORMANT features
does not produce a clear AUC improvement over MFCC in isolation, and
combining MFCC with COVAREP alone is numerically \emph{worse}
($0.682$) than MFCC alone. This is consistent with our earlier
decision (Section~\ref{sec:features}) not to claim COVAREP/FORMANT as
a performance contribution, and instead describe their inclusion as a
feature-engineering choice. One genuinely interesting exploratory
finding is that FORMANT features alone---just 5 dimensions per
frame---reach AUC~$0.688\pm0.060$, comparable to the 120-dim MFCC
representation; this is consistent with the clinical literature
linking formant structure (vowel space) to depressed speech
\cite{scherer2016vowel}, and is worth follow-up with a larger, dedicated
seed count rather than being treated as confirmed here.

% ============================================================
\section{Discussion}
\label{sec:discussion}

\subsection{A Missing-Data Sentinel in the OpenFace AU Stream}
\label{sec:sentinel}

\begin{figure}[t]
  \centering
  \includegraphics[width=\columnwidth]{fig12_sentinel_fix.png}
  \caption{Effect of the OpenFace $-100$ missing-data sentinel. Bars
    show the original 5-seed protocol ($\pm1$~SD); black points are
    individual seeds and grey lines connect the same seed across arms.
    Removing sentinel-carrying frames before normalization lifts the
    face-only configuration from chance to AUC~$0.708$ at this 5-seed
    protocol; the annotated significance test is instead the scaled,
    correctly-computed $n{=}30$ re-test (Section~\ref{sec:sentinel}).
    The tri-modal configuration is essentially unchanged, indicating
    the recovered facial signal is largely redundant with audio and
    text.}
  \label{fig:sentinel}
\end{figure}

Our ablation reports the face-only configuration at chance
(AUC~$=0.513\pm0.113$, Table~\ref{tab:ablation}). Rather than accept this
as evidence that CLNF Action Units carry no depression signal, we audited
the feature stream itself and found a data-handling fault.

The six binary AU channels (\texttt{AU04\_c}, \texttt{AU12\_c},
\texttt{AU15\_c}, \texttt{AU23\_c}, \texttt{AU28\_c}, \texttt{AU45\_c})
encode tracking failure with a $-100$ sentinel rather than a valid
$\{0,1\}$ presence value. Our preprocessing---following common practice
on this corpus---filters frames only by \texttt{confidence}~$\geq0.5$,
which does \emph{not} remove these rows: every participant inspected
retains sentinel-contaminated frames (mean~$3.3\%$ of retained frames on
train, up to $22.9\%$ for individual participants on dev).

The consequence is mechanical rather than subtle. Because the pipeline
then z-scores each channel per participant, a value two orders of
magnitude outside the real $\{0,1\}$ range dominates both moments. For
participant~308's \texttt{AU12\_c} channel ($9.0\%$ sentinels), the
channel statistics become $\mu=-8.856$, $\sigma=28.664$ instead of
$\mu=0.157$, $\sigma=0.364$; after normalization the genuine frames
occupy a span of $0.035$ rather than $2.749$---a $\approx79\times$
compression that leaves the encoder a near-constant input on 6 of its 20
channels.

Dropping frames that carry the sentinel (a sentinel marks a failed
track, so the whole frame is untrustworthy) and changing nothing else
raises the face-only configuration from AUC~$0.513\pm0.113$ to
$0.708\pm0.042$ at an initial 5-seed pilot---every seed
improves, and seed variance falls $2.7\times$ (Fig.~\ref{fig:sentinel}).
An exact two-sided Wilcoxon signed-rank test cannot reach $p<0.05$ at
$n{=}5$ under any outcome (minimum $p{=}0.0625$; the same constraint
noted for the multi-task comparison in Section~\ref{sec:ablation}), so
this pilot was underpowered by construction. We therefore re-ran both
arms (original sentinel-affected features vs.\ cleaned) from scratch
at $n{=}30$ seeds: 28/30 seeds improve, AUC~$0.506\pm0.131\to
0.678\pm0.048$, Wilcoxon signed-rank $p{=}9.3{\times}10^{-6}$,
Cohen's~$d{=}1.28$. We report this $n{=}30$ result as the paper's
confirmatory evidence for the sentinel effect. This is a data-quality
effect rather than a modelling one.

Applying the same fix to the single-task tri-modal configuration
(\texttt{ConfigurableFusionModel}, the shared architecture behind every
row of Table~\ref{tab:ablation}), however, changes almost nothing
(AUC~$0.697\pm0.051 \to 0.702\pm0.057$; 2/5~seeds improve; $p{=}0.313$).
We additionally re-trained the proposed multi-task model itself
(\texttt{MultiTaskFusionModel}---a distinct architecture with its own
prediction heads and modality-dropout regularizer), 5 seeds under each
face-feature version: AUC changes by $-0.009$ ($0.698\pm0.033 \to
0.689\pm0.040$; 1/5~seeds improve; $p{=}0.125$) and F1 by $+0.030$
($0.609\pm0.050 \to 0.639\pm0.051$; 2/5~seeds improve; $p{=}0.500$),
neither significant. We read the three tested results together as
follows: the facial signal destroyed by the sentinel is real and
recoverable, but it is largely \emph{redundant} with what audio and
text already supply, on both architectures---consistent with
Section~\ref{sec:robustness}, where zero-masking the face channel at
test time costs almost nothing on the multi-task model. Two
consequences follow. First, this paper's headline tri-modal
(single-task) \emph{and} proposed multi-task numbers are both
confirmed unaffected by the fault by direct re-test, not merely
assumed. Second, any DAIC-WOZ study that reports a weak or
chance-level \emph{visual} branch should verify sentinel handling
before concluding that facial behaviour is uninformative; we cannot
tell from published descriptions how widely this is checked.

We separately tested the alternative hypothesis that per-participant
z-scoring itself removes the clinically-relevant signal, by computing
85~interpretable dynamics features (per-AU intensity mean, temporal
standard deviation, robust range, frame-to-frame change rate, activation
fraction, and presence/transition rates) directly from \emph{raw},
unnormalized AU intensities, and testing depressed vs.\ non-depressed
differences on the pooled train$+$dev set ($N{=}141$) with
Mann-Whitney~$U$ and Benjamini-Hochberg FDR control. No feature survives
correction (smallest raw $p{=}0.115$), and the held-out test split was
therefore never queried. In this corpus we find no measurable
blunted-affect signal in absolute AU intensity or variability---the
sentinel, not the normalization scheme, accounts for the face branch's
failure.

\subsection{Interviewer-Prompt Shortcut Audit}
\label{sec:shortcut}

\begin{figure}[t]
  \centering
  \includegraphics[width=\columnwidth]{fig9_shortcut_audit.png}
  \caption{Interviewer-prompt shortcut audit. The same
    \texttt{ConfigurableFusionModel} text configuration, training
    schedule, and 5~seeds are used in all three conditions; only which
    speaker's transcript turns are retained differs, and all conditions
    are evaluated on an identical matched participant set. An
    Ellie-only model---which never sees a single word spoken by the
    patient---achieves the highest AUC on both splits, indicating the
    label leaks through interview structure rather than through
    depression-specific patient language. Error bars: $\pm1$~SD across
    seeds.}
  \label{fig:shortcut}
\end{figure}

A growing body of work questions the validity of DAIC-WOZ text results.
Burdisso et al.\ \cite{burdisso2024therapist} show that text models
including the virtual interviewer's (Ellie's) utterances can reach
F1~=~0.90---not by detecting depression, but by exploiting \emph{which}
questions Ellie asks: the semi-structured interview branches into a
mental-health-history line of questioning for some participants, so the
interviewer's own words leak the label. An independent large-scale
reproducibility audit \cite{patapati2025invalid} similarly finds that
most published DAIC-WOZ classifiers fail basic task-specificity checks.

We audit this shortcut directly inside our own pipeline rather than
taking the concern on faith. Using the identical
\texttt{ConfigurableFusionModel}, training schedule, and 5 seeds as the
Section~\ref{sec:ablation} ablation, we train the text-only
configuration under three conditions that differ \emph{only} in which
speaker's transcript turns are retained, with the TF-IDF vectorizer
refit per condition on train only and all three conditions evaluated
on an identical, leakage-free participant set (matched $N$: train~=~107,
dev~=~33, test~=~46):

\begin{itemize}
  \item \textbf{Participant-only} (used throughout this paper):
  test AUC~=~0.610, F1~=~0.456.
  \item \textbf{Participant~+~Ellie}: test AUC~=~0.681, F1~=~0.516.
  \item \textbf{Ellie-only}---the model never sees a single word spoken
  by the patient: test AUC~=~\textbf{0.738}, F1~=~\textbf{0.568}
  (dev AUC~=~0.841).
\end{itemize}

\noindent Fig.~\ref{fig:shortcut} summarizes these three conditions on
both splits.

The participant-only condition closely matches the text-modality
number reported in Table~\ref{tab:ablation} (dev
AUC~$0.5714\pm0.0397$ here vs.\ $0.5614\pm0.0431$ there, across the
same 5 seeds); the small difference is due to the matched-$N$
restriction (dev $N{=}33$ vs.\ $34$), which excludes one participant
whose transcript lacks usable text under one of the three conditions
so that all three conditions are evaluated on an identical
participant set. This confirms the participant-only condition here
and in the main ablation are the same model under the same protocol,
differing only in this minor, documented sample restriction. The
Ellie-only model---with
zero access to patient language---outperforms the participant-only
model by a wide margin on both dev and test. Inspecting the logistic-
regression coefficients of an auxiliary classifier trained on Ellie-only
TF-IDF features shows the terms most predictive of the depressed class
are follow-up questions Ellie asks only of participants who have
already disclosed treatment history: \emph{``seek help,'' ``how long ago
were you diagnosed,'' ``still go to therapy now.''} The terms most
predictive of not-depressed are neutral small talk (\emph{``awesome,''
``great thanks,'' ``laughter''}). The label is leaking through interview
\emph{structure}, not through any depression-specific language produced
by the participant.

This finding directly supports our participant-only design choice
(Section~\ref{sec:features}) and reframes our comparatively modest
headline numbers (F1~=~0.585, AUC~=~0.706, Table~\ref{tab:sota}) as
evidence of validity rather than a weakness relative to literature
values obtained with interviewer text included. We report this
audit as a methodological contribution independent of the model's
absolute performance: a positive result here (participant-only text
also inflated) would have equally been reported, since the audit's
purpose is to characterize the corpus, not to flatter this paper's
numbers.

\subsection{Limitations}
\label{subsec:limitations}

\textbf{Small dev set ($N{=}34$):} The official dev split provides
insufficient statistical power to establish a significant AUC
difference between tri-modal and audio-only configurations
(paired bootstrap, 2,000 resamples, $p=0.859$)---audio-only is
numerically higher (Table~\ref{tab:ablation}), but the gap is not
distinguishable from sampling noise at this sample size. We
additionally report test-split results ($N{=}47$): F1\,$=0.585$,
AUC\,$=0.706$ (see Table~\ref{tab:sota}). The test AUC exceeds
the dev AUC, providing evidence of generalization. The F1 decrease
($0.629\to0.585$) is attributable to threshold mismatch between
the train-tuned threshold (Section~\ref{subsec:threshold}) and the
test distribution. Pooled 10-fold
cross-validation over all 188 participants
(Section~\ref{sec:kfold}) corroborates that a single fixed split is
an unreliable estimator at this sample size: F1 and TPR-gap both
vary substantially across folds.

\textbf{Audio-only AUC exceeds tri-modal AUC:} Table~\ref{tab:ablation}
shows that audio-only achieves AUC~$=0.721\pm0.013$ vs.\
the proposed tri-modal model AUC~$=0.658\pm0.029$. This is a known phenomenon
in multimodal learning on small datasets: adding modalities introduces
additional noise dimensions under limited supervision, which can
\emph{hurt} AUC (discrimination) while improving F1 (calibrated
threshold decision) through multi-task regularization. Three
observations support this interpretation: (1)~audio-only has the
lowest AUC variance ($\pm0.013$), indicating it learns the most
stable discriminative boundary; (2)~tri-modal achieves higher F1
($0.629$ vs.\ $0.593$) by combining weaker signals with optimized
threshold; (3)~leave-one-modality-out ablation on the dev split shows
audio removal causes the largest AUC drop ($-0.03$), consistent with
audio carrying the primary discriminative signal \emph{on dev}---though
Section~\ref{sec:robustness}'s test-time masking experiment finds this
does not hold on the held-out test split, where removing audio does
not hurt and removing text is instead the critical failure mode. On a
larger dataset, the additional modalities are expected to provide
complementary signal rather than noise. This gap is not an artefact of
the face-branch sentinel fault described in Section~\ref{sec:sentinel}:
repairing that fault on the single-task tri-modal configuration
(\texttt{ConfigurableFusionModel}, the architecture common to every row
of Table~\ref{tab:ablation}) leaves its AUC within noise of the
original value ($0.697\pm0.051\to0.702\pm0.057$, $p{=}0.313$). We
separately re-trained the proposed multi-task model
(\texttt{MultiTaskFusionModel}, a distinct architecture with its own
heads and modality dropout) under the corrected face features, 5 seeds
each arm: AUC changes by $-0.009$ ($p{=}0.125$) and F1 by $+0.030$
($p{=}0.500$), neither significant. We therefore treat the
audio-only/tri-modal gap reported here as reflecting genuine
small-sample multimodal noise rather than a fixable data error, now
confirmed by direct re-test on both architectures rather than assumed
from the single-task result alone. We flag one remaining caveat for
consistency with Section~\ref{sec:ablation}: both AUC values compared
in this paragraph (audio-only $0.721\pm0.013$ and the proposed
multi-task model $0.658\pm0.029$) still come from the original 5-seed
sample, not the scaled 70-seed test. That scaled test targeted F1 on
the single-task-vs-multi-task comparison specifically and found the
multi-task model's F1 stable across seeds ($0.657\pm0.037$ over
$n{=}70$), which is reassuring for the multi-task side of the AUC
comparison here; it did not, however, rescale the audio-only
configuration or the AUC metric itself, so we cannot rule out that
the audio-only AUC estimate shifts at a larger seed count the way the
single-task \emph{F1} baseline did. We report the 5-seed AUC gap as
observed rather than as confirmed at scale, and flag scaling the
audio-only configuration to a larger seed count as future work,
consistent with how we treat this same open question elsewhere in the
paper.

\textbf{PHQ-8 as ground truth:} PHQ-8 is a self-report instrument,
not a clinician-administered diagnosis. Label noise inherent in
self-report limits the ceiling performance of any supervised model;
AUC~=~0.70 may represent the empirical ceiling on this label source.

\textbf{CPU training:} All reported results use CPU training on the
official splits. BERT-based text encoding and larger-scale GPU training
are avenues for future improvement; no specific performance gain
is claimed here as these remain untested.

\textbf{Binary gender labels:} DAIC-WOZ provides binary gender only.
Fairness analysis should be extended to intersectional groups and
non-binary identities as datasets with richer demographic annotation
become available.

\begin{figure}[t]
  \centering
  \includegraphics[width=\columnwidth]{fig10_cross_corpus.png}
  \caption{Zero-shot cross-corpus generalization in two modalities.
    Each DAIC-WOZ-trained model is applied without fine-tuning to an
    independent, Chinese-language corpus. Both collapse to chance.
    The audio probe's wide interval reflects EATD-Corpus's small
    sample ($N{=}162$); the text probe's near-invisible interval
    reflects WU3D's $200\times$ larger sample ($N{=}32{,}540$), making
    that null result statistically decisive. Error bars: 95\%
    bootstrap CI, 2,000 resamples.}
  \label{fig:crosscorpus}
\end{figure}

\textbf{Cross-corpus generalization:} All results above are evaluated
within DAIC-WOZ (same corpus, recording setup, interview format, and
population), so they characterize generalization to unseen
\emph{participants} but not to unseen \emph{corpora}. To probe this
directly, we take the audio-only configuration---the strongest and
most stable unimodal model (Table~\ref{tab:ablation}, dev
AUC~$=0.721\pm0.013$; re-verified on a fixed-seed checkpoint at dev
AUC~$=0.735$)---and evaluate it \emph{zero-shot}, with no
fine-tuning, on EATD-Corpus \cite{shen2022eatd}: 162 Chinese-language
volunteers responding to free-form emotional prompts, labeled
depressed/non-depressed via the Self-Rating Depression Scale
(SDS~$\geq$~53; 30 depressed, 132 non-depressed). Both corpora are
reduced to the same 120-dimensional MFCC+$\Delta$+$\Delta^2$
representation before inference. Performance collapses from
AUC~$=0.735$ in-domain to AUC~$=0.531$ (95\% CI $[0.406, 0.653]$,
2,000 bootstrap resamples)---statistically indistinguishable from
chance (AUC~$=0.5$)---and F1 drops from $0.56$ to $0.30$. The
confusion matrix shows the model recovers only 14 of 30 depressed
participants (recall~$=0.47$) while producing 49 false positives
among 132 non-depressed participants. This indicates the acoustic
representation learned is corpus- and language-specific---plausibly
reflecting DAIC-WOZ's controlled Wizard-of-Oz interview acoustics and
English prosody---rather than a transferable, language-agnostic
depression biomarker. We report this as a direct empirical bound on
clinical deployability: held-out same-corpus evaluation, however
rigorous, does not by itself support claims of generalization to a
new population, language, or recording setting.

We replicate this pattern in the text modality with a second,
independent, and substantially larger cross-corpus probe. A
participant-only text-only checkpoint (identical
\texttt{ConfigurableFusionModel} architecture; five seeds trained and
each verified on the DAIC-WOZ held-out test split before selection,
since Table~\ref{tab:ablation}-level dev performance alone was not
sufficient to guarantee test-time generalization---this is the same
caution the audio probe above illustrates; best-generalizing seed:
test AUC~$=0.684$) is evaluated zero-shot, through the same
DAIC-WOZ-fit TF-IDF vectorizer, on WU3D \cite{wang2020wu3d}: 32,540
Sina Weibo users (10,324 depressed, 22,216 normal), whose labels were
rechecked against DSM-5 criteria by professional annotators rather
than inferred from subreddit self-selection. This is a complete
language shift (English clinical transcript~$\to$~Chinese
social-media text) layered on top of the domain shift already tested
with EATD-Corpus. Performance collapses to AUC~$=0.506$ (95\% CI
$[0.5008, 0.5121]$, 2,000 bootstrap resamples)---an order of
magnitude tighter interval than the EATD probe, owing to WU3D's
$200\times$ larger sample, making the null result statistically
decisive rather than merely plausible. The mechanism is directly
observable in the features themselves: 55.9\% of WU3D users produce
an all-zero TF-IDF vector under the English-fit vectorizer (mean 4.11
non-zero features out of 1,000), confirming the collapse reflects
near-total vocabulary mismatch rather than a subtler representational
failure. Two independent cross-corpus probes, in different
modalities, on different corpora, in different languages, both
collapse from strong in-domain performance to chance
(Fig.~\ref{fig:crosscorpus})---consistent
evidence that this pipeline (and, we suspect, comparable
small-corpus DAIC-WOZ systems in general) learns corpus-specific
rather than universal depression markers. We were unable to run an
equivalent probe for the face modality: the only depression-labeled
corpus we identified with OpenFace-extracted Action Units in a format
compatible with our pipeline, CMDC \cite{zou2023cmdc}, requires an
EULA-gated access process incompatible with this study's timeline. We
note explicitly that this gap became \emph{more} rather than less
important once the sentinel fault was corrected: before the fix, a
chance-level face branch would have made a cross-corpus face probe
uninformative, but the corrected branch reaches AUC~$0.708\pm0.042$
in-domain (Section~\ref{sec:sentinel}) and therefore has something that
could meaningfully fail to transfer. We flag this as the clearest
outstanding gap in the paper's cross-corpus evidence rather than
explaining it away.

\textbf{Diagnostic specificity: depression, or distress in general?}
DAIC-WOZ carries one self-report scale and no second diagnosis, so
within-corpus accuracy cannot separate a depression detector from a
general-distress detector---the central criticism of
\cite{patapati2025invalid}. We test it directly. The participant-only
text model is scored, without fine-tuning, on a publicly available
English social-media corpus labeled by mental-health category
\cite{sarkar2022mentalhealth}, and we ask whether its
depression score can tell those categories apart. Because the corpus's
\emph{Normal} class has a median length of 10~words against 87--137 for
every distress class, unmatched contrasts would largely measure post
length (and hence how many TF-IDF terms fire); we therefore pair each
positive post with an unused negative post of similar word count
($\pm25\%$) and report the matched figure. Labels here derive from the
community a post was drawn from, not from clinical assessment, and are
treated as category indicators rather than diagnoses.

The model transfers to this domain (Depression vs.\ Normal, matched
AUC~$=0.678$, 95\% CI $[0.656, 0.701]$)---unlike the WU3D probe above,
because the English vocabulary survives (mean 53 non-zero TF-IDF terms
per document vs.\ 4.11 for WU3D). But it also fires on other distress
categories (Suicidal vs.\ Normal $0.569$; Anxiety vs.\ Normal $0.581$;
Stress vs.\ Normal $0.550$), and its ability to \emph{separate}
depression from those categories is weak: Depression vs.\ Suicidal
$0.573$ $[0.557, 0.589]$, Depression vs.\ Anxiety $0.568$
$[0.552, 0.584]$, Depression vs.\ Stress $0.626$ $[0.611, 0.641]$.
These sit above chance---the model is not purely a distress
detector---but far below its depression-vs-normal performance. We read
this as partial support for \cite{patapati2025invalid}: much of what the
model detects is shared across distress conditions, with a smaller
depression-specific component on top---the same qualified conclusion
Section~\ref{sec:symptom} reaches from the complementary angle of
within-corpus symptom-level heterogeneity.

One consequence deserves stating plainly, since it concerns safety
rather than performance. The model separates depression from
suicidality at AUC~$0.573$---barely above chance. This is expected: the
PHQ-8 label it was trained on is precisely PHQ-9 \emph{minus} the
suicidality item, so the model has never been exposed to a suicidality
label, and nothing in this work supports its use for suicide-risk
assessment. We report the contrast only to bound what the model
distinguishes, not to propose a capability; a system that cannot
separate these two conditions must not be deployed where that
distinction carries clinical weight.

\textbf{Learning curve---would more data actually help?} The paper
repeatedly attributes non-significant effects to insufficient sample
size without demonstrating that more data would help. We test this
directly: the tri-modal \texttt{ConfigurableFusionModel} (identical
architecture and training schedule as the main ablation) is trained on
stratified subsamples of the 107-participant train pool at 30\%, 50\%,
70\%, 85\%, and 100\% (3 seeds each), evaluated throughout on the
fixed, never-subsampled dev set ($N{=}34$). AUC rises from
$0.572\pm0.036$ at $N{=}32$ to $0.725\pm0.031$ at $N{=}91$
(Fig.~\ref{fig:learningcurve}); the small dip at the full pool
($N{=}107$: AUC~$=0.693\pm0.025$) falls within the overlapping
standard-deviation bands of the $N{=}91$ point and is consistent with
seed noise rather than a genuine reversal. The curve has not visibly
plateaued within the available data, supporting the paper's repeated
claim that this pipeline is data-limited rather than
architecture-limited---a claim tested here rather than merely
asserted.

\begin{figure}[t]
  \centering
  \includegraphics[width=\columnwidth]{fig11_learning_curve.png}
  \caption{Learning curve: tri-modal AUC/F1 on the fixed dev split as
    the training pool is subsampled from the full 107 participants
    (3~seeds per point, stratified by class). AUC increases with
    $N$ and shows no clear plateau, indicating this pipeline remains
    data-limited at the corpus's current size. Error bars: $\pm1$~SD.}
  \label{fig:learningcurve}
\end{figure}

\subsection{Clinical Significance}

This work addresses three open problems in computational psychiatry.
First, multimodal coverage: the ablation study's error analysis finds
5 of 34 dev participants misclassified by the best single unimodal
configuration (audio-only) but correctly classified by tri-modal
fusion, and only 1 participant misclassified by every configuration
tested. This case-level result is consistent with the multi-channel
nature of clinical observation, even though tri-modal's \emph{aggregate}
AUC is lower than audio-only's (Section~\ref{sec:discussion})---the two
findings describe different things: tri-modal trades some aggregate
discrimination for coverage of cases audio-only misses. We report
this as a specific, checkable observation rather than a general claim
that fusion "confirms" superiority. Second, training-time fairness: the
Equalized Odds loss is integrated into the objective function---not
applied post-hoc---making fairness a first-class design constraint,
though our 10-seed evaluation shows it is insufficient for significant
bias reduction at $N{=}34$, motivating larger-scale replication.
Third, interpretability: three independent XAI methods converge on
audio as the most informative dev-split modality, and gradient
saliency places the clinically-established facial marker AU04 within
its top~5 attributed channels (though not ranked first), providing a
partial evidence chain a clinician can audit rather than a fully
clean one.

% ============================================================
\section{Conclusion}
\label{sec:conclusion}

We presented a systematic validity audit of a tri-modal depression-risk
pipeline on DAIC-WOZ, surfacing three failure modes---a silent
data-corruption sentinel in the visual stream, an interview-structure
shortcut in the text stream, and non-transferable, corpus-specific
representations in audio and text---that plausibly explain a
substantial share of the wide performance variance reported across
the DAIC-WOZ literature. Within the same controlled testbed, the most
promising design-choice result we observed is symptom-level PHQ-8
multi-task learning: a controlled ablation
(7~configurations\,$\times$\,5~seeds\,=\,35~runs, bootstrap~CI,
paired bootstrap significance) shows the multi-task model raises
mean macro-F1 from $0.607\pm0.063$ to $0.629\pm0.021$ with markedly
reduced seed variance, and accuracy~$=68.8\%$ on the DAIC-WOZ dev
split in the original 5-seed sample. A power analysis of that 5-seed
paired comparison indicated it was underpowered
($n{\approx}67$ needed for 80\% power), so we scaled the paired
Wilcoxon test to $n{=}70$ seeds (Section~\ref{sec:ablation}). This
confirms a significant effect (51/70~seed pairs favor
multi-task, $p{=}8.1{\times}10^{-7}$, $d{=}0.72$, larger than the
$d{=}0.34$ the original power calculation assumed), while the
single-task baseline's 70-seed mean ($0.597\pm0.076$) is close to the
originally-reported $0.607\pm0.063$---the original 5-seed test's main
limitation was that $n{=}5$ is structurally too small for an exact
Wilcoxon test to reach significance, not a badly biased mean.
The combined multi-task intervention (prediction heads plus extended
199-dim audio features, Section~\ref{sec:ablation}) is therefore now
reported as a statistically confirmed result rather than an
encouraging trend, though it confirms that combined intervention
rather than isolating multi-task learning alone as the causal factor.
Two additional design choices---cross-modal
attention fusion and Equalized Odds fairness training---remain
evaluated at $n{=}10$ paired Wilcoxon seeds and are reported as
exploratory, not headline, contributions.

One of the two confirmed significant effects in this paper is not a
modelling result at all but a data-quality one: OpenFace's $-100$
tracking-failure sentinel survives the standard confidence filter and,
under per-participant z-scoring, compresses the genuine AU signal
$\approx$79$\times$, holding the face-only configuration at chance.
Removing the affected frames raises face-only AUC from
$0.513\pm0.113$ to $0.708\pm0.042$ at the original 5-seed protocol; a
correctly-computed, scaled re-test at $n{=}30$ seeds confirms this as
real (28/30 seeds improving, Wilcoxon $p{=}9.3{\times}10^{-6}$,
$d{=}1.28$)---the original 5-seed report's $p{=}0.031$ could not be
reproduced and is corrected in Section~\ref{sec:sentinel}. Repairing
the sentinel leaves both the single-task
tri-modal result ($p{=}0.313$) and the proposed multi-task result
($p{=}0.125$ AUC, $p{=}0.500$ F1, both re-tested with 5 seeds)
unchanged---so both headline numbers are confirmed by direct re-test,
not merely assumed---but any DAIC-WOZ study reporting a weak visual
branch should check sentinel handling before concluding that facial
behaviour is uninformative
(Section~\ref{sec:sentinel}).

The remaining two design-choice comparisons are honest rather than
favorable: (1)~cross-modal attention shows a positive trend over
concatenation (F1\,$+0.024$, AUC wins 6/10~seeds) but does not reach
statistical significance at $N{=}34$ ($p{=}0.496$); (2)~the Equalized
Odds loss does not significantly reduce the gender TPR gap at this
sample size ($p{=}0.813$, wins 3/10~seeds). We report both as
negative/inconclusive findings that characterize the limits of these
techniques on small clinical cohorts, rather than reframing them as
successes. By contrast, the multi-task comparison---rescaled to
$n{=}70$ paired seeds once a power analysis flagged the original
$n{=}5$ test as underpowered---does reach significance
($p{=}8.1{\times}10^{-7}$, $d{=}0.72$); we highlight this as an
example of what resolving an underpowered comparison can reveal, and
recommend the same rescaling wherever feasible for the two
comparisons that remain inconclusive.
Test-split evaluation ($N{=}47$) yields F1\,$=0.585$ and
AUC\,$=0.706$, exceeding the dev AUC and providing evidence of
generalization beyond the dev set. The test-split TPR gap is
nominally~0.000, but pooled 10-fold cross-validation across all
188 participants (Section~\ref{sec:kfold}) shows this single-split
result is not stable (mean TPR gap $0.415\pm0.211$ across folds),
underscoring that fixed-split fairness metrics on $N{<}50$ cohorts
should be corroborated with resampling before being reported as
findings. Test-time missing-modality masking
(Section~\ref{sec:robustness}) shows the modality-dropout regularizer
gives partial, not uniform, robustness: face is dispensable---and
Section~\ref{sec:sentinel} establishes this is redundancy rather than a
broken channel, since repairing the sentinel fault lifts face in
isolation without moving the tri-modal result---but
removing text collapses test AUC from~0.706 to~0.435. Symptom-level
breakdown (Section~\ref{sec:symptom}) finds accuracy varies
meaningfully across the eight PHQ-8 items (presence-absence
macro-F1~$0.209$--$0.424$) but is highest on the most generic
affective symptoms, offering only partial evidence against the
concern that DAIC-WOZ classifiers learn disorder-general rather than
disorder-specific cues. An internal audio-feature ablation
(Section~\ref{sec:audio_ablation}) finds MFCC alone matches the full
199-dim MFCC+COVAREP+FORMANT combination (AUC $0.710$ vs.\ $0.711$),
supporting our decision not to claim COVAREP/FORMANT as a performance
contribution, though FORMANT alone is a notable 5-dimensional
exploratory result. Calibration analysis
(Section~\ref{sec:calibration}) finds moderate probability
miscalibration (ECE $0.102$/$0.157$ on dev/test) and a
false-negative-skewed dev error profile, both relevant to any future
deployment decision. Three convergent explainability methods
(attention rollout, gradient$\times$input saliency, occlusion)
complete the framework. Future work includes GPU training, BERT text
encoding, probability recalibration, and replication on larger
same-corpus extensions such as E-DAIC to determine whether the
multi-task, attention, and fairness effects become statistically
significant with more data. Closing the cross-corpus gap identified
above will likely require domain adaptation or multi-corpus training
rather than scale alone, and further cross-corpus probes (e.g.\ on
in-the-wild multimodal corpora such as LMVD) are needed to map the
boundary of this limitation across modalities beyond audio.

% ============================================================
\section*{Ethics Statement}
This study performs secondary analysis of the DAIC-WOZ corpus, which
was collected by the University of Southern California Institute for
Creative Technologies (USC ICT) under its own institutional human-
subjects protocol, with documented participant informed consent
\cite{gratch2014}. The authors did not collect any new data from human
participants and accessed only de-identified interview recordings and
questionnaire scores under the standard USC ICT data-use agreement;
no additional institutional review was required for this secondary
analysis at the authors' institution.

% ============================================================
\section*{Data Availability Statement}
The DAIC-WOZ corpus is not publicly redistributable; it is available
to qualified researchers directly from USC ICT
(\url{https://dcapswoz.ict.usc.edu/}) upon signing a data-use
agreement. EATD-Corpus, WU3D, and the Kaggle mental-health text corpus
\cite{sarkar2022mentalhealth}, used for the cross-corpus and
diagnostic-specificity probes (Section~\ref{sec:discussion}), are all
publicly available datasets. All source code, trained model
configurations, and result artifacts (excluding the restricted raw
dataset) are available at
\url{https://github.com/DevMursLab/Thesis}, with the exact version
underlying this manuscript archived at
\url{https://doi.org/10.5281/zenodo.22806159} (v1.0.0).

\subsection*{Reproducibility Details}
All experiments use Python~3.10, PyTorch~$\geq$2.0, scikit-learn~$\geq$1.3,
and librosa~$\geq$0.10 (exact pinned versions in the repository's
\texttt{requirements.txt}). The 5-seed ablation and multi-task
experiments use seeds $\{42, 1, 7, 123, 2024\}$; the 10-seed fairness
and attention-vs-concatenation experiments use seeds
$\{42,1,7,13,21,100,2024,7777,555,88\}$; the scaled multi-task
significance test (Section~\ref{sec:ablation}) re-trains both the
single-task and multi-task arms from scratch for seeds $1$--$70$
inclusive. \texttt{torch.manual\_seed}
and \texttt{numpy.random.seed} are set identically per run; exact
bitwise reproducibility across GPU/CPU hardware is not guaranteed
due to non-deterministic CUDA kernels, though result variation across
reruns at fixed seed was observed to be negligible relative to the
inter-seed variance reported throughout this paper.

% ============================================================
\section*{Author Contributions}
Md.~Mursalin: conceptualization, methodology, software, formal
analysis, investigation, data curation, writing---original draft,
visualization. Md.~Shovon (Supervisor): validation, writing---review
and editing, supervision. Both authors read and approved the final
manuscript.

% ============================================================
\section*{Conflict of Interest Statement}
The authors declare that they have no known competing financial
interests or personal relationships that could have appeared to
influence the work reported in this paper.

% ============================================================
\section*{Funding Statement}
This research received no external grant funding from any
funding agency in the public, commercial, or not-for-profit
sectors.

% ============================================================
\section*{Acknowledgment}
The DAIC-WOZ corpus was provided by the University of Southern
California Institute for Creative Technologies under a restricted-data
agreement. EATD-Corpus, WU3D, and the Kaggle mental-health text corpus
\cite{sarkar2022mentalhealth} are publicly available datasets used for
cross-corpus and diagnostic-specificity evaluation.

\subsection*{Declaration of Generative AI and AI-Assisted Technologies}
During manuscript preparation, the corresponding author used an
AI-assisted coding and writing tool (Claude, Anthropic) to help draft
and revise experiment scripts, generate figures from computed results,
identify inconsistencies during proofreading, and edit prose for
clarity. All experimental results, statistical analyses, and
scientific claims were designed, executed, and verified by the
authors, who reviewed and take full responsibility for all AI-assisted
content in this manuscript.

% ============================================================
\begin{thebibliography}{20}

\bibitem{who2023}
World Health Organization, ``Depressive disorder (depression),''
\textit{WHO Fact Sheets}, Geneva, Switzerland, 2023.
[Online]. Available: \url{https://www.who.int/news-room/fact-sheets/detail/depression}

\bibitem{gratch2014}
J.~Gratch, R.~Artstein, G.~Lucas, G.~Stratou, S.~Scherer,
A.~Nazarian, R.~Wood, J.~Boberg, D.~DeVault, S.~Marsella,
D.~Traum, A.~Rizzo, and L.-P.~Morency,
``The Distress Analysis Interview Corpus of human and computer
interviews,'' in \textit{Proc. Int. Conf. Lang. Resources Eval.
(LREC)}, Reykjavik, Iceland, 2014, pp.~3123--3128.

\bibitem{ringeval2017}
F.~Ringeval, B.~Schuller, M.~Valstar, J.~Gratch, R.~Cowie,
S.~Scherer, S.~Mozgai, N.~Cummins, M.~Schmitt, and M.~Pantic,
``AVEC 2017: Real-life depression, and affect recognition workshop
and challenge,'' in \textit{Proc. 7th ACM Int. Workshop Audio/
Visual Emotion Challenge (AVEC)}, Mountain View, CA, USA, 2017,
pp.~3--9.

\bibitem{gong2017}
Y.~Gong and C.~Poellabauer,
``Topic modeling based multi-modal depression severity assessment,''
in \textit{Proc. 7th ACM Int. Workshop Audio/Visual Emotion
Challenge (AVEC)}, 2017, pp.~53--60.

\bibitem{williamson2016}
J.~R.~Williamson, E.~Godoy, M.~Cha, A.~Schwarzentruber,
P.~Khorrami, Y.~Gwon, H.~T.~Kung, C.~K.~Dagli, and T.~F.~Quatieri,
``Detecting depression using vocal, facial and semantic
communication cues,'' in \textit{Proc. 6th Int. Workshop
Audio/Visual Emotion Challenge (AVEC)}, 2016, pp.~11--18.

\bibitem{ray2019}
A.~Ray, S.~Kumar, R.~Reddy, P.~Mukherjee, and R.~Garg,
``Multi-level attention network using text, audio and video
for depression prediction,'' in \textit{Proc. 9th Int. Workshop
Audio/Visual Emotion Challenge (AVEC)}, 2019, pp.~81--88.

\bibitem{burdisso2024therapist}
S.~Burdisso, E.~Reyes-Ram\'irez, E.~Villatoro-Tello,
F.~S\'anchez-Vega, A.~Lopez-Monroy, and P.~Motlicek,
``DAIC-WOZ: On the validity of using the therapist's prompts
in automatic depression detection from clinical interviews,''
in \textit{Proc. 6th Clinical Natural Language Processing
Workshop}, 2024, pp.~82--90.

\bibitem{patapati2025invalid}
S.~V.~Patapati, I.~Pendyala, M.~Ambati, P.~Kunadharaju,
P.~Kokati, \emph{et al.},
``Most DAIC-WOZ depression classifiers are invalid, they don't
learn task-specific features: Preliminary findings from a
large-scale reproducibility study,'' in \textit{Companion Proc.
27th Int. Conf. Multimodal Interaction (ICMI)}, 2025.

\bibitem{wang2020wu3d}
Y.~Wang, Z.~Wang, C.~Li, Y.~Zhang, and H.~Wang,
``A multitask deep learning approach for user depression
detection on Sina Weibo,''
\textit{arXiv preprint arXiv:2008.11708}, 2020.

\bibitem{zou2023cmdc}
B.~Zou, J.~Han, Y.~Wang, R.~Liu, S.~Zhao, L.~Feng, X.~Lyu, and H.~Ma,
``Semi-structural interview-based Chinese multimodal depression
corpus towards automatic preliminary screening of depressive
disorders,''
\textit{IEEE Trans. Affective Comput.}, vol.~14, no.~4,
pp.~2823--2838, 2023.

\bibitem{sarkar2022mentalhealth}
S.~Sarkar,
``Sentiment analysis for mental health,''
Kaggle dataset, 2022. [Online]. Available:
\url{https://www.kaggle.com/datasets/suchintikasarkar/sentiment-analysis-for-mental-health}

\bibitem{covarep2014}
G.~Degottex, J.~Kane, T.~Drugman, T.~Raitio, and S.~Scherer,
``COVAREP---A collaborative voice analysis repository for speech
technologies,'' in \textit{Proc. IEEE Int. Conf. Acoust. Speech
Signal Process. (ICASSP)}, Florence, Italy, 2014, pp.~960--964.

\bibitem{scherer2013voice}
S.~Scherer, G.~Stratou, J.~Gratch, and L.-P.~Morency,
``Investigating voice quality as a speaker-independent indicator
of depression and PTSD,'' in \textit{Proc. Interspeech},
Lyon, France, 2013, pp.~847--851.

\bibitem{scherer2016vowel}
S.~Scherer, G.~Lucas, J.~Gratch, A.~Rizzo, and L.-P.~Morency,
``Self-reported symptoms of depression and PTSD are associated with
reduced vowel space in screening interviews,''
\textit{IEEE Trans. Affective Comput.}, vol.~7, no.~1,
pp.~59--73, 2016.

\bibitem{hardt2016}
M.~Hardt, E.~Price, and N.~Srebro,
``Equality of opportunity in supervised learning,'' in
\textit{Proc. Adv. Neural Inf. Process. Syst. (NeurIPS)},
Barcelona, Spain, 2016, pp.~3315--3323.

\bibitem{barocas2019}
S.~Barocas, M.~Hardt, and A.~Narayanan,
\textit{Fairness and Machine Learning: Limitations and Opportunities}.
2019. [Online]. Available: \url{https://fairmlbook.org}

\bibitem{ekman1978}
P.~Ekman and W.~V.~Friesen,
\textit{Facial Action Coding System (FACS): A Technique for the
Measurement of Facial Action}.
Palo Alto, CA, USA: Consulting Psychologists Press, 1978.

\bibitem{abnar2020}
S.~Abnar and W.~Zuidema,
``Quantifying attention flow in transformers,'' in
\textit{Proc. 58th Annu. Meeting Assoc. Comput. Linguistics
(ACL)}, Online, 2020, pp.~4190--4197.

\bibitem{shrikumar2017}
A.~Shrikumar, P.~Greenside, and A.~Kundaje,
``Learning important features through propagating activation
differences,'' in \textit{Proc. 34th Int. Conf. Mach. Learn.
(ICML)}, Sydney, Australia, 2017, pp.~3145--3153.

\bibitem{shen2022eatd}
Y.~Shen, H.~Yang, and L.~Lin,
``Automatic depression detection: An emotional audio-textual
corpus and a GRU/BiLSTM-based model,''
in \textit{Proc. IEEE Int. Conf. Acoust., Speech, Signal Process.
(ICASSP)}, Singapore, 2022, pp.~6247--6251.

\bibitem{bouthillier2021}
X.~Bouthillier, P.~Delaunay, M.~Bronzi, \emph{et al.},
``Accounting for variance in machine learning benchmarks,''
in \textit{Proc. Mach. Learn. Syst. (MLSys)}, 2021.

\bibitem{tsai2019multimodal}
Y.-H.~H.~Tsai, S.~Bai, P.~P.~Liang, J.~Z.~Kolter,
L.-P.~Morency, and R.~Salakhutdinov,
``Multimodal transformer for unaligned multimodal language
sequences,'' in \textit{Proc. 57th Annu. Meeting Assoc.
Comput. Linguistics (ACL)}, Florence, Italy, 2019,
pp.~6558--6569.

\bibitem{baltrusaitis2016}
T.~Baltru\v{s}aitis, P.~Robinson, and L.-P.~Morency,
``OpenFace: An open source facial behavior analysis toolkit,''
in \textit{Proc. IEEE Winter Conf. Appl. Comput. Vis. (WACV)},
Lake Placid, NY, USA, 2016, pp.~1--10.

\bibitem{danylenko2026}
I.~Danylenko and O.~Unold,
``Common pitfalls and recommendations for use of machine
learning in depression severity estimation: DAIC-WOZ study,''
\textit{Appl. Sci.}, vol.~16, no.~1, p.~422, 2026.

\end{thebibliography}

\end{document}
